\documentclass[lettersize,journal]{IEEEtran}
\usepackage{amsmath,amsfonts}
\usepackage{array}
\usepackage{textcomp}
\usepackage{stfloats}
\usepackage{url}
\usepackage{verbatim}
\usepackage{graphicx}
\hyphenation{op-tical net-works semi-conduc-tor IEEE-Xplore}
\ifCLASSOPTIONcompsoc
\usepackage[caption=false, font=normalsize, labelfont=sf, textfont=sf]{subfig}
\else
\usepackage[caption=false, font=footnotesize]{subfig}
\fi


\usepackage[linesnumbered, ruled]{algorithm2e}%要求显示行号和算法标准框
\SetKwRepeat{Do}{do}{while}% repeat..until可以用do...while来实现和显示
\SetKwInOut{Init}{Initialization}


\usepackage{hyperref}
% \hypersetup{colorlinks=true}
\usepackage{booktabs}
\usepackage{multirow}
\usepackage[table]{xcolor}
\usepackage{makecell}
\usepackage{color}
\usepackage{amssymb}
\usepackage[numbers,sort&compress]{natbib}

\newtheorem{theorem}{Theorem}[section]
\newtheorem{lemma}{Lemma}[section]
\newtheorem{corollary}{Corollary}[section]
\newtheorem{proof}{Proof}[section]
% \captionsetup[figure]{name={Fig.},singlelinecheck=off}

\begin{document}
%
% paper title
% Titles are generally capitalized except for words such as a, an, and, as,
% at, but, by, for, in, nor, of, on, or, the, to and up, which are usually
% not capitalized unless they are the first or last word of the title.
% Linebreaks \\ can be used within to get better formatting as desired.
% Do not put math or special symbols in the title.
\title{Data-Driven Chance-Constrained Multiple-Choice Knapsack Problem:\\Model, Algorithms, and Applications}
%
%
% author names and IEEE memberships
% note positions of commas and nonbreaking spaces ( ~ ) LaTeX will not break
% a structure at a ~ so this keeps an author's name from being broken across
% two lines.
% use \thanks{} to gain access to the first footnote area
% a separate \thanks must be used for each paragraph as LaTeX2e's \thanks
% was not built to handle multiple paragraphs
%

\author{Xuanfeng~Li,~\IEEEmembership{Student Member,~IEEE,}
        Shengcai~Liu,~\IEEEmembership{Member,~IEEE,}
        Jin~Wang,
        Xiao~Chen,
        Yew-Soon~Ong,~\IEEEmembership{Fellow,~IEEE,}
        Ke~Tang,~\IEEEmembership{Fellow,~IEEE}
\thanks{Xuanfeng Li and Ke Tang are with the Guangdong Provincial Key Laboratory of Brain-Inspired Intelligent Computation, Department of
Computer Science and Engineering, Southern University of Science and Technology, Shenzhen 518055, China (lixf2020@mail.sustech.edu.cn, tangk3@sustech.edu.cn).}
\thanks{Shengcai Liu is with the Centre for Frontier AI Research (CFAR), Agency for Science, Technology and Research (A*STAR), Singapore (Liu\_Shengcai@cfar.a-star.edu.sg).}
\thanks{Jin Wang and Xiao Chen are with Huawei Reliability Lab (wangjin171@huawei.com, chenxiao18@huawei.com).}
\thanks{Yew-Soon Ong is with the School of Computer Science and Engineering, Data Science and Artificial Intelligence Research Centre, Nanyang Technological University (NTU), and also with the Centre for Frontier AI Research (CFAR), Agency for Science, Technology and Research (A*STAR), Singapore (asysong@ntu.edu.sg).}
}

\maketitle
\begin{abstract}
The multiple-choice knapsack problem (MCKP) is a classic NP-hard combinatorial optimization problem. 
Motivated by several significant practical applications, this work investigates a novel variant of MCKP called data-driven chance-constrained multiple-choice knapsack problem (DDCCMCKP), where the item weight is a random variable with unknown probability distribution.
We first present the problem formulation of DDCCMCKP, and then establish two benchmark sets.
The first set contains synthetic instances, and the second set is devised to simulate a real-world application scenario of a certain telecommunication company.
To solve DDCCMCKP, we propose a data-driven adaptive local search (DDALS) algorithm.
The main merit of DDALS lies in evaluating solutions with chance constraints by data-driven methods, under the condition of unknown distributions and only historical sample data being available.
The experimental results demonstrate the effectiveness of the proposed algorithm and show that it is superior to other baselines. 
Additionally, ablation experiments  confirm the necessity of each component in the algorithm.
Our proposed algorithm can serve as the baseline for future research, and the code and benchmark sets will be open-sourced to further promote research on this challenging problem.
\end{abstract}

% Note that keywords are not normally used for peerreview papers.
\begin{IEEEkeywords}
Combinatorial optimization, data-driven approach, multiple-choice knapsack problem, evolutionary algorithm.
\end{IEEEkeywords}

\IEEEpeerreviewmaketitle

\section{Introduction}

% 背包问题是什么，目标是什么，变量是什么。并给出应用
% 虽然在过去围绕如何求解MCKP已经有了许多研究工作，但是在实际应用中会出现一种情况，即物品的重量并不是确定的值而是随机变量。

\IEEEPARstart{T}{he} multiple-choice knapsack problem (MCKP) is a variant of the classic knapsack problem \cite{sinha1979multiple}.
MCKP finds practical applications in various areas, such as supply chain optimization \cite{sharkey2011class}, target market selection \cite{taaffe2008target}, investments portfolio selection \cite{nauss19780}, transportation programming \cite{zhong2010multiple}, and component selection in the information technology systems \cite{kwong2010optimization}.
Specifically, MCKP involves a set of items, each possessing a weight and a value. The items are divided into several disjoint classes and only one item can be selected in each class \cite{sinha1979multiple,kellerer2004multiple}. The objective of MCKP is to select exactly one item from each class such as to maximize the total value while adhering to a predefined weight limit (i.e., knapsack constraint).
Over the past several decades, plenty of methods for solving MCKP have been proposed, including greedy methods \cite{zemel1980linear}, branch-and-bound \cite{dyer1984branch}, dynamic programming \cite{dudzinski1987exact}, approximation algorithms \cite{gens1998approximate}, and heuristic search \cite{hifi2006reactive}.

% Despite the huge research progress achieved on solving MCKP, a scenario one can often encounter in practice is that the weights of the items are not fixed values but random variables.
% In such cases, the classical knapsack constraint is no longer sufficient to model the problem.
Despite of significant research success in addressing the MCKP, those advanced methods are incapable of handling problems in real-world scenarios where item weights are randomly distributed instead of fixed.
To handle random variables, a common approach is to model deterministic constraints as chance constraints, i.e., the probability (confidence level) of total weight (random variable) of the selected items not exceeding a given deterministic value is no less than a certain threshold \cite{doerr2020optimization,shapiro2021lectures}.
Furthermore, the specific probability distributions of the random variables could be unknown \cite{ji2021data,xie2021distributionally}.
Instead, one only has historical sample data to characterize the random variables. 
The collection of sample data for each item is assumed to be independent and uncorrelated.
This variant of MCKP is referred to as data-driven chance-constrained multiple-choice knapsack problem (DDCCMCKP). 


% 首先介绍第一个具体例子：端到端 5G网络配置，每一种实现方式会带来不同的cost和time delay，其中cost是确定值，而time delay是未知的随机变量
A typical application of DDCCMCKP is the configuration of fifth-generation (5G) end-to-end (E2E) network.
In the area of 5G for business \cite{nasrallah2018ultra}, user plane time delay is one of the key performance indicators for low latency communication \cite{5g2017view,parvez2018survey}.
%latency is considered one of the key performance indicators.
As illustrated in Figure~\ref{fig:5G-e2e}, the complete packet transmission path in a 5G cellular network consists of several segments, i.e., eNB, backhaul, core network, Internet/cloud, and the transport between core network and Internet/cloud.
More specifically, these segments correspond to classes, where for each class there exist different options (items) and only one option can be selected to implement the class.
The total time delay of a packet transmission can be calculated as the sum of the delays induced by all the selected options, i.e., $(T_{eNB}+T_{Backhaul}+T_{Core}+T_{Transport}+T_{Internet})*2$.
Moreover,  the packet transmission delays of the implementation options are random variables with unknown distributions as they are affected by various factors, including workloads, network configurations, broadcast environment, and random failures of software and hardware.
The configuration of 5G E2E network is to select an implementation option for each segment.
Different configurations have different total costs and time delays, where the costs are fixed values, but time delays are unknown random variables.
The objective is to minimize the total cost of network configuration while ensuring that the probability of meeting the delay is less than a given threshold under a certain value.

Another application of DDCCMCKP is supply chain planning and optimization in modern industries \cite{sharkey2011class, taaffe2008target}.
The supply-chain operations reference (SCOR) model, published by the Supply-Chain Council \cite{huan2004review}, is a useful tool for the upper management of a firm to design and reconfigure the supply chain to achieve better performance.
In SCOR, the E2E supply chain consists of four distinct processes: source, make, deliver, and plan (see Figure \ref{fig:cpo-mi}). Specifically, these processes can be further subdivided into supplier selection, R\&D testing, production, transportation, and recycling handling, which correspond to different classes along the supply chain. Therefore, each class has various configuration options that can be influenced by random factors, such as market demand, cost fluctuations, and technology configurations.

\begin{figure}[!t]
\centering
\includegraphics[width=3in]{chap1_5G_E2E-formal.pdf}
\caption{5G E2E network service model \cite{parvez2018survey}} 
\label{fig:5G-e2e}
\end{figure}

\begin{figure}[!t]
\centering
\includegraphics[width=3in]{chap1_SCOR_E2E_formal.pdf}
\caption{SCOR model-based supply chain infrastructure \cite{huan2004review}} 
\label{fig:cpo-mi}
\end{figure}

% Despite that DDCCMCKP has not been tackled before, it is worth investigating because it has broad and significant applications as mentioned above. 
Since DDCCMCKP has not been paid research attention, it merits thorough investigation due to its extensive and noteworthy applications.
Compared to MCKP, DDCCMCKP faces greater challenges due to the complexity of its constraints. 
There are two main difficulties: it is a chance-constrained optimization problem, and the distributions of the random variables are unknown.
Therefore, a data-driven approach that does not rely on specific distribution information but rather available sample data to evaluate solutions must be proposed.

The main contributions are summarized as follows:

\begin{enumerate}
\item 
From the problem perspective, we define DDCCMCKP from several real-world scenarios and formulate the model. 
To the best of our knowledge, this is the first endeavor to address this specific problem. 
To further promote research on this problem, we have designed two sets of benchmarks. 
The first set is a simulated benchmark with a simple mechanism, while the other set is to simulate a real-world business requirement of a telecommunication company.
The generated benchmarks will be made publicly available to promote further research.
%online for everyone to follow and test the algorithms.

\item 
From the algorithmic perspective, we propose a novel algorithm named data-driven adaptive local Search (DDALS). 
The first novelty of DDALS is that we evaluate solutions with chance constraints by a data-driven method. 
The second is that we design surrogate weights using statistics of sample data to ensure the performance of the search algorithm. 
Furthermore, our algorithm can adaptively jump between small and larger search regions for efficient searching of high-quality solutions. 
%adopts the framework of local search algorithm, which 
%can adaptively jump between small and larger regions for efficient searching of solutions. 
Lastly, we improve the proportion of feasible solutions through carefully designed filtering mechanisms.

\item From the computational perspective, we experimentally confirm the effectiveness of DDALS.
We adapt several commonly-used approaches as the baselines, including greedy search (Greedy), genetic algorithm (GA), and estimation of distribution algorithm (EDA). 
Experimental results show that DDALS outperforms other baseline methods.
The proposed algorithm can be served as a baseline for comparison experiments on this challenging problem.
%for future research, and the code and benchmark set will be open-sourced to encourage further studies.
\end{enumerate}

The rest of the paper is organized as follows.
In Section \ref{section:rw}, we present a comprehensive review of related works. 
Section \ref{section:np} provides a formal definition of DDCCMCKP. 
In Section \ref{section:ddals}, DDALS is introduced and a detailed explanation of its implementation is provided. 
In Section \ref{section:bd}, the benchmarks we designed are described in details. 
The experimental results and analyses are given in Section \ref{section:cs}. 
Finally, we summarize our findings and provide suggestions for future research in Section \ref{section:con}.

\section{Related works}
\label{section:rw}
In this section, we will review the relevant literature. We first review existing research on the traditional MCKP. The second category of research focuses on chance-constrained optimization. Since the methods developed in robust optimization can solve chance-constrained optimization by constructing proper uncertainty sets \cite{han2016robust,joung2020robust}, we further review robust optimization and distributionally robust optimization. While robust optimization approach minimizes the worst-case scenario under the uncertainty of random variables, the goal of distributionally robust optimization is to ensure the robustness of the optimal solution in the distribution-unknown situations.

\subsection{Algorithms for MCKP}
There have been various stochastic variations of the classical 0-1 MCKP that have been extensively studied in the literature. 
For instance, Dyer et al.\cite{dyer1984branch} proposed a branch and bound method to solve the MCKP. 
Pisinger \cite{pisinger1995minimal} presented a minimal algorithm including linear relaxation, greedy algorithm and dynamic programming. 
Lamanna et al. \cite{lamanna2022two} introduced a new algorithm for solving the multidimensional MCKP based on Kernel search, which divided the solution process of the algorithm into two phases with complementary objectives. 
Gens and Levner \cite{gens1998approximate} proposed a fast approximation algorithm for the MCKP with a 4/5 approximation ratio of the optimal solution using binary search. 
He et al. \cite{he2016improved} developed an improved approximate binary search algorithm with an approximation ratio of $3+(1/2)^t$. 
Hifi et al. \cite{hifi2006reactive} proposed a reactive local search algorithm for the multidimensional MCKP. 
Mkaouar et al.  \cite{mkaouar2020solving} introduced a method based on the artificial bee colony algorithm to solve the multidimensional MCKP. 
In general, existing researches on MCKP cannot be applied, since the solutions in DDCCMCKP can not be evaluated by deterministic approaches.

\subsection{Chance-Constrained Optimization}
Chance-constrained optimization plays a critical role in stochastic optimization \cite{delage2010percentile,doerr2020optimization,shapiro2021lectures}. 
Kleinberg et al. \cite{kleinberg1997allocating} focused on packing multiple bursty connections with varying transmission rates on a link and considered stochastic knapsack problems with known Bernoulli distribution. 
Goel and Indyk \cite{goel1999stochastic} studied the knapsack problem when the jobs had stochastic processing requirements or sizes. They proposed polynomial time approximation scheme (PTAS) when the jobs were all known to follow either Poisson or exponential distributions. 
Goyal and Ravi\cite{goyal2010ptas} proposed a PTAS that used a parametric linear programming reformulation and yielded near-optimal profit efficiently. The random variables were assumed to follow a known Gaussian distribution. 
Shabtai et al. \cite{shabtai2018relaxed} concentrated on the chance-constrained knapsack problem with known normal distributions. They presented a fast PTAS for the problem that run in poly($n/\epsilon$) time. 
Zhang et al. \cite{8637163} studied trajectory prediction with chance constraint, assuming that  future position of the moving obstacle could be represented as Gaussian distributions based on the predicted statistics.
Chai et al. \cite{9269349} studied trajectory optimization with chance constraint with known probability density function of random variables defined on the measurable set.
Yang and Chakraborty \cite{yang2018algorithm} exploited the geometry of the non-convex feasible region of the chance-constrained knapsack problem in the variance-mean plane and transformed the problem into a sequence of deterministic knapsack problems to find an optimal solution.
All these works mainly focused on theoretical analysis of situations that the distribution of random variables is known, which are not applicable to our problem as well.

\subsection{Robust Optimization}
\label{subsection:ro}
A robust optimization approach aims to construct an uncertainty set $\mathcal{U}_{\epsilon}\subseteq \mathcal{U}$ such that a feasible solution of the form $g(x,\xi)\leq 0, \forall \xi \in \mathcal{U}_{\epsilon}$ will also be feasible with probability at least $1-\epsilon$ with respect to $P$, i.e., $\mathcal{P}(g(x,\xi)\leq 0)\geq 1-\epsilon, \forall \xi \in \mathcal{U}$, which has the same form with chance constraint. 
Han \cite{han2016robust} considered a chance-constrained knapsack problem where each item had an independent normally distributed random weight with a known mean and variance, and proposed an efficient pseudo-polynomial time algorithm based on the robust optimization approach. 
Joung and Lee\cite{joung2020robust} proposed a robust optimization-based heuristic algorithm where the optimal solution could be obtained by solving the ordinary binary knapsack problem iteratively. 

Distributionally robust optimization was proposed as an improvement of robust optimization when the distribution of the random variable was unknown and only observable through a finite training dataset.
Mohajerin Esfahani and Kuhn \cite{mohajerin2018data} showed that distributionally robust optimization problems over Wasserstein balls could be reformulated as finite convex programs. 
Ji and Lejeune \cite{ji2021data} studied distributionally robust chance-constrained programming optimization problems with data-driven Wasserstein ambiguity sets. 
The proposed algorithmic and reformulation framework were applied to all types of distributionally robust chance-constrained optimization problems, including the multidimensional knapsack problem.
Additionally, Xie \cite{xie2021distributionally} proposed that the distributionally robust chance-constrained problem could be reformulated as a conditional value-at-risk constrained optimization problem. 
The problem with binary decision variables admitted a big-M free formulation and could be solved by a branch and cut algorithm. 
Gao and Kelywegt \cite{gao2022distributionally} derived a strong duality reformulation of the corresponding distributionally robust stochastic optimization problem and construct approximate worst-case distributions explicitly. 
Note that most of these work can not be applied. 
We have tried to make modifications to Ji’s \cite{ji2021data} and Yang's \cite{yang2018algorithm} algorithm to make it applicable to our problem, but the results show that both of them exhibited unsatisfactory performance. (see Section \ref{subsection:jy} for details).

\section{Notations and problem definition}
\label{section:np}
We give the formulation of DDCCMCKP in Eq. \ref{formu_DDCCMCKP}. 
Consider $m$ mutually disjoint classes $N_1,\ldots,N_m$ of items. 
Each item $j\in N_i$ has an associated cost $c_{ij}$ and weight $w_{ij}$, where $c_{ij}, j\in N_i, \forall i\in \mathcal{M}=\{1,\ldots,m\}$ are fixed values while $w_{ij},j\in N_i,\forall i\in \mathcal{M}$ are unknown random variables. 
The weight of item $j\in N_i$ is characterized by $L$ sample data, denoted as $d_{ijl},i\in \{1,\ldots,m\},j\in N_i, l\in \{1,\ldots,L\}$.
Note that the sample data of each item is independent and identically distributed (i.i.d.).
The goal is to select one item from each class to minimize the total cost, subject to keeping the total weight below the knapsack’s capacity $W$ with probability higher than the given confidence level $P_0$.
By introducing the binary variables $x_{ij}$ which take on value 1 if the item $j$ is chosen in class $N_i$, the problem can be formulated as:

\begin{subequations}
    \begin{align}
    & & \min \quad       &\sum\limits_{i=1}^m\sum\limits_{j\in N_i} c_{ij} x_{ij}    \\
    & & \text{s.t.}\quad &\mathcal{P}(\sum\limits_{i=1}^m\sum\limits_{j\in N_i} w_{ij} x_{ij}\leq W)\geq P_0 \label{cc}\\
    & &                  &\sum\limits_{j\in N_i}x_{ij}=1, \forall i \in \mathcal{M}\\
    & &                  &x_{ij}\in \{0,1\},\forall i \in \mathcal{M}, j\in N_i 
    \end{align}
    \label{formu_DDCCMCKP}
\end{subequations}

 To avoid trivial solutions, we assume that the following inequality holds:
\begin{equation}
    \sum\limits_{i=1}^m \min_{1\leq l \leq L, 1\leq j\leq n_i}d_{ijl} \leq W\leq \sum\limits_{i=1}^m \max_{1\leq l \leq L, 1\leq j\leq n_i}d_{ijl}
\end{equation}

Note that MCKP is NP-hard \cite{sinha1979multiple}. 
Supposing that all samples of each item's weight $w_{ij}$ are the same, then DDCCMCKP reduces to traditional MCKP. 
Therefore, MCKP can be seen as a special case of DDCCMCKP, which indicates that DDCCMCKP is also NP-hard. 
It motivates us to adopt an algorithm based on a heuristic search framework to solve DDCCMCKP.

\section{Data-driven adaptive local search}
\label{section:ddals}
As mentioned in Section \ref{section:rw}, although there have been many works in the fields of MCKP, chance-constrained optimization and robust optimization, none of the methods could be directly applied to our problem. 
In order to tackle the DDCCMCKP, we present a novel approach called data-driven adaptive local search (DDALS) in this work. 
Specifically, DDALS consists of three main parts, confidence level evaluation, surrogate weight design and local search framework. 
Finally, we will analyze the complexity of the algorithm.

\subsection{Confidence Level Evaluation}
Since the sample data is assumed to be i.i.d., it is necessary to make full use of the information contained in the sample data, since we have no information about the true distribution. 
Given a solution $S$, a natural idea is to calculate all possible combinations of the data and count the proportion of these combinations that meet the requirements. 
This proportion is an estimate of the real confidence level $P$ of the solution. 
However, obviously, a brute-force method is not applicable since it requires a calculation complexity of $L^m$, which is related to the problem size $m$ and sample size $L$. 
It may be feasible for small-scale problems, but it becomes impractical for large-scale problems. 
Therefore, based on the above considerations, we propose an exact evaluation method for small-scale problems that is much faster than brute-force enumeration. 
For large-scale problems, we propose an accelerated Monte Carlo evaluation method.

\subsubsection{Exact evaluation}
Exact evaluation is based on a simple idea: counting and heaps \cite{ferreira1996fast}.
Assume that the sample data of items has been already arranged from largest to smallest. 
Solution $S = [s_1,s_2,\ldots,s_m]$ is under evaluated. 
The details of exact evaluation are listed as follows:
\begin{itemize}
    \item Maintain a maximum heap queue $Q$ with the length $m(1-P_0)L^m$. Every time $Q$ pops the largest sum in the current heap queue \cite{ferreira1996fast};
    \item Initialization: push $Sum_1 = d_{s_1,1}+\ldots+d_{s_m,1}$ into $Q$;
    \item Iteration: $Q$ pops the largest sum in $D$:
    \begin{equation*}
        \begin{aligned}
           D =  &[d_{s_1l_1},d_{s_2l_2},\ldots,d_{s_ml_m}].\\
        \end{aligned}
    \end{equation*}
    Then push neighbors of the popped $D$ whose sum are smaller than that of $D$ into $Q$, listed as follows:
    \begin{equation*}
        \begin{aligned}
        &[d_{s_1l_1+1},d_{s_2l_2},\ldots,d_{s_ml_m}],\\
        &[d_{s_1l_1},d_{s_2l_2+1},\ldots,d_{s_ml_m}],\\
        &\ldots,\\
        &[d_{s_1l_1},d_{s_2l_2},\ldots,d_{s_ml_m+1}]\\
    \end{aligned}
    \end{equation*}
    Note that duplicates are not allowed in $Q$.
    \item Repeat until the pop sum of $D$ is smaller than $W$ or the iteration times $k$ reach $ \lceil(1-P_0)L^m\rceil$.
\end{itemize}

We have the following theorem,
\begin{theorem}
The popped sum in the $k^{th}$ iteration is exactly the $k^{th}$ largest sum.
\end{theorem}

Detailed proofs can be found in Appendix 2. If the ratio of the number to the total number of combinations is greater than $1-P_0$, then we determine that the solution is infeasible.

\subsubsection{Accelerated Monte Carlo evaluation}
Monte Carlo simulation is a method that uses random numbers to solve many computational problems \cite{metropolis1949monte}. 
In this work, we focus on situations where the knapsack constraints $W$ are tight at a high confidence level $P_0$, which means feasible solutions account for a very small proportion of the entire solution space of the problem. 
This observation motivates us that the evaluation procedure can be greatly accelerated if most of the infeasible solutions can be quickly eliminated. 
Specifically, assuming that the sample data of the items has been sorted from largest to smallest, we can check the $l^{th}$ largest data combination $d_{1l}, d_{2l},...,d_{ml}$ for the combination of items to be solved, i.e.,$D_l = \sum_{i=1}^{m}d_{il}$, by comparing with the weight threshold $W$ to directly filter out most of the infeasible solutions. 
Therefore, the evaluations times required during the algorithm operation can be significantly reduced. 
The process can be summarized as follows.
\begin{itemize}
    \item Construct a filtered list according to the minimum prime factorization for $K=L^m\times (1-P_0 )$.
    \item Only a limited size of combinations will be selected to construct the filtered list $R\ll m!$;
    \item Calculate the sum of the data of each item of the screening list $R$, and check whether there is at least one combination satisfying $\sum_{i=1}^{m}d_{il} \geq W$. If not, the solution will be evaluated by Monte Carlo simulation.
\end{itemize}

\subsection{Surrogate Weight Design}
It is not feasible to directly obtain the cost-effectiveness of each item in a data-driven scenario, which impedes the algorithm’s ability to search for better solutions. 
To overcome this challenge, we construct a surrogate cost-effectiveness based on the statistic of sample data. 
Specifically, we calculate the sample mean and standard deviation of the item weights, $\mu_{ij}=\mathbb{E}(w_{ij}),\sigma^2_{ij}=Var(w_{ij}),i\in \{1,...,m\}, j \in N_i$. 
Then item weights are represented by the weighted sum of statistical characteristics,
\begin{equation}\label{Definition: surrogate weights}
    \tilde{w}_{ij}:=\mu_{ij}+\lambda\sigma _{ij}, \forall{i}=1,...,m, j\in N_i.
\end{equation}
where $\lambda $ is the weighted parameter. 
This approach provides a statistical measure for guiding the search process towards optimal solutions.

\subsection{Local Search Framework}
Based on the previously data-driven solution evaluation method and surrogate weight design, we now introduce the overall algorithmic framework of DDALS. 
Local search is a widely used heuristic search framework for efficiently solving various combinatorial optimization problems \cite{hifi2006reactive,9877844}. 
Thus, our algorithm adopts the framework of local search algorithm, which can adaptively jump between small and large regions, achieving efficient search for solutions. 
The block diagram of DDALS shown in Figure \ref{fig:Flow_Diagram} . 
DDALS is composed of the following five components: constructive procedure (CP) for solution initialization, a local swap search (LSS) operator for improving the current solution, a degradation (Degrade) operator for escaping local optima, and a further swap search (FSS) operator for searching in a larger feasible region. 
Since limited sample data cannot exactly reflect the true distribution, we additionally incorporate a solution feasibility enhancement (SFE) to improve the feasibility of the output.

\begin{figure}[!t]
\centering
\includegraphics[width=3in]{chap4_FlowGramLS.pdf}
\caption{The block diagram of DDALS} 
\label{fig:Flow_Diagram}
\end{figure}

\begin{algorithm}
    \SetAlgoVlined
    \caption{Data-Driven Adaptive Local Search}
    \label{alg:frameworkofDDALS}
    \KwIn{A DDCCMCKP benchmark; Maximum iteration times, $MaxIter$; weighting factor, $\lambda$}
    \KwOut{the optimal feasible solution $S_{opt}$}
    $S_{cp} \gets ConstructProcedure(\lambda)$, cost list, MC list\;
    \While{$t\leq MaxIter$}{
        $S_{t}\gets LocalSwapSearch(S_t)$\;
        \If{$S_t$ is a better solution}{$S^* \gets S_{t}$}
        $S_{t}\gets Degrade(S_t)$\;
        $t=t+1$\;
    }
    $S_{opt}\gets FurtherSwapSearch(S^*)$\;
    Update $S_{opt}$ through solution feasibility enhancement\;
    Return $S_{opt}$\;
\end{algorithm}

Algorithm \ref{alg:frameworkofDDALS} illustrates the framework of DDALS. 
The inputs of DDALS include the benchmark, maximum iteration times $MaxIter$, and a user-defined weighting factor $\lambda$. 
We first initialize the solution $S_{cp}$ using CP (line 1), then iteratively improve the quality of solutions. 
At iteration $t$, LSS attempts to obtain a better feasible solution by operating on $S_t$ (line 3). 
If LSS finds an improvement, the optimal solution $S^*$ is updated (line 4-5). 
The current solution $S_t$ will be degraded by randomly selecting an item for replacement to escape a local optimum (line 6). 
After iterations end, FSS is employed on $S_{opt}$ to search for a better solution (line 8). 
The final solution $S_{opt}$ will be output by SFE. 

The detailed designs of the components are list as follows:

\subsubsection{Constructive procedure}
The initialization procedure is based on greedy strategy \cite{kellerer2004multiple}.
Utilizing the item weights defined in Definition \ref{Definition: surrogate weights}, we calculate the cost-effective variables $u_{ij}=v_{ij}/\tilde{w}_{ij}$ for each item $j \in N_i$, then sort the items in descending order of $u_{ij}$ for each class $N_i$. 
$S_{cp}$ is constructed by selecting the item with the largest $u_{ij}$ in each class. 
If the current $S_{t}$ is infeasible,  the item which has the largest $\tilde{w}_{i_0j}$ in $S_t$ is replaced by an item with a smaller weight in the same class $i_0$. 
The procedure iterates until a feasible $S_{cp}$ is obtained.
If all items in class $i_0$ are tried but still infeasible,$S_{cp}$ will be constructed by the smallest $\tilde{w}_{ij}$ in each class.

\subsubsection{Local swap search}
LSS applies a swapping strategy for picked items to improve the current solution $S_t$. 
If the obtained solution is feasible, each swap is authorized. 
LSS will try all the items with lower cost in each class $N_i$ and selects the item with the lowest cost from all authorized swaps. 

\subsubsection{Degrade}
Degrade randomly replaces an item $j_{i_0}$ of class $i_0$ in the current solution $S_t$. 
If an infeasible solution is generated, it randomly selects a lighter item in $i_0$. 
If there is no suitable item in the $i_0$, it randomly switches to another class and tries again until a new feasible solution $S_t$ is obtained.

\subsubsection{Further swap search}
FSS iterates through all class combinations $(i_0,i_1)$ of the current optimal solution $S_t$ to find a better and feasible solution. 
If it finds a pair of items with a smaller cost, $S_t$ will be updated. The algorithm ends when all pairs of items have been searched.

\subsubsection{Solution feasibility enhancement}
In practice, due to the demand for high confidence level $P_0$, we tend to minimize the cost of the solution while ensuring its feasibility, rather than blindly pursuing the minimization of the cost, leading to solutions that do not satisfy constraints, which renders them useless in practice. 
It is a crucial requirement in many practical applications.
However, it is difficult to guarantee the real feasibility of the optimal solution due to limited sample size. 
Note that the sample size to achieve smaller errors for arbitrarily distributed random variables will be very large (See Appendix 3 for theoretical analysis between estimation error and sample size). 
Therefore, we design three filtering variants to enhance the feasibility of output solution, while not increasing the cost too much. 
In the iterations of DDALS algorithm, two list of feasible solutions are recorded, each with a capacity of 30. 
The first records 30 feasible solutions with the lowest cost encountered so far, named Cost List.
The second records 30 solutions with the largest confidence level encountered so far, named MC List.

\begin{itemize}
    \item Original: Output the optimal solution of FSS. 
    \item Variant 1: The minimum-cost solution with a higher confidence level than $P_0$. For example, it should satisfy $P \geq 0.995$ if $P_0 = 0.99$. 
    \item Variant 2: All solutions will be scored according to the order in the list they are in. 
    Each candidate solution has two scores, which will be added together. 
    Choose the solution with the highest score as the final output.
    \item Variant 3: Take the first 10 solution (if the amount is larger than 10) in Cost List as the final output and choose one manually.
\end{itemize}

\subsection{Complexity Analysis}

\subsubsection{Exact evaluation}
First, if the solution can be directly determined infeasible in the fast screening phase, the time complexity takes $O(1)$. 
Otherwise, in order to verify whether it is feasible, the maximum number of iterations that the algorithm needs to perform is $T=L^m (1-P_0)-l_{min}^m$, where $l_{min}$ is the index that the combination of $l_{min}$ samples of solution is larger than $T_{\max}$. 
The maximum heap size is $k_0+(m-1)T=\sum_{i=1}^m \binom{m}{i} (l_{min}+1)^{m-i}+(m-1) L^m (1-P_0 ) -(m-1) l_ {min}^m=O(mT)$.
Since at each iteration it needs to sort the heap once, the total number of times is $O(\log(mT))=O(\log m+\log T)=O(\log T),m\ll L$.
Thus, the total time complexity of the filtered algorithm is $O(mL\log L+\sum_{i=1}^m \binom{m}{i}(l_{min}+1)^{m-i}+ T\log T)= O(\max\{T\log T,ml_{min}^{m-1}, m^2 l_{min}^{m-2},...,ml_{min} \})$.

\subsubsection{Accelerated Monte Carlo evaluation}
The complexity of the accelerated Monte Carlo evaluation is $O((|G|+\alpha \times MonteCarlo)\times Eval)$, where $MonteCarlo$ denotes the set number of Monte Carlo simulations and $|G|$ denotes the size of fast screening combination. 
If the solution is an infeasible solution determined by fast screening, $\alpha =0$, otherwise $\alpha=1$.

\subsubsection{DDALS}
\begin{itemize}
    \item 
    CP includes sorting and solution evaluation. Up to $m \times N$ times of evaluation are required. 
    The complexity is $O(Eval\times mN)$, where $N=\max|N_i|$.
    \item 
    Similar to CP, LSS require up to $m \times N$ times of evaluation to search a possible better solution.
    The complexity is $O(Eval\times mN)$.
    \item 
    The worst case is that only one combination of items is feasible in the whole solution space. 
    The complexity of Degrade is $O(Eval \times mN)$.
    \item 
    FSS iterates through a maximum of combinations of classes $\binom{m}{2}\times N^2$.
    The complexity of FSS is $O(Eval \times m(m-1) N^2/2) = O(Eval \times m^2 N^2 )$.
\end{itemize}
Note that SFH only takes $O(1)$. Therefore, the total time complexity of DDALS is $O(Eval\times mN + MaxIter\times (Eval\times mN+Eval\times mN)+Eval\times N^2 m^2 +1) = O(Eval\times \max\{N^2 m^2,MaxIter\times mN\})$. 

\section{Benchmark Design}
\label{section:bd}
Since DDCCMCKP is a new problem that has not been solved before, it is necessary to design benchmarks to evaluate the performance of algorithms in different scenarios. 
Therefore, we designed two sets of benchmarks to cover a wide range of problem sizes and complexities, each of which consists of multiple benchmarks of different scales, in order to comprehensively evaluate the algorithm’s performance in different scenarios.

\begin{figure*}[htpb]
\centering
\subfloat[LAB]{\includegraphics[width=2.0\columnwidth]{chap6_pdf_of_different_distribution_lab.png}%
\label{fig:pdf_LAB}}
\hfil
\subfloat[APP]{\includegraphics[width=2.0\columnwidth]{chap6_pdf_of_different_distribution_app.png}%
\label{fig:pdf_APP}}
\caption{Probability density function of items. The horizontal axis represents the time delay and the vertical axis represents the probability density.}
\label{fig: APPLAB}
\end{figure*}

The first set of benchmarks aims to simulate the case where the delay of items follows different common continuous distributions, named LAB.
The distribution of each item is randomly chosen from the set of uniform, truncated normal, fatigue-life, bimodal and gamma distributions. 
The mean of items is uniformly distributed in $U(2,8)$, i.e., $\mu\sim U(2,8)$, and the variance is uniformly distributed in $U(1,19)$, i.e., $\sigma^2 \sim U(1,19)$, where $U(2,8)$ denotes a uniform distribution over the interval [2,8]. 
In the case of the gamma distribution, the mean value is $\mu\sim U(0.5,2.5)$ and variance $\sigma^2\sim U(0.05,0.625)$. 
The cost, representing the value of taking the item, is generated randomly between $U(1,10)$, corresponding to uncorrelated situation in classic MCKP benchmark\cite{mkaouar2020solving}, which is a widely used benchmark in the literature. 
The curve of probability density function of different distributions of the items involved in LAB is shown in Figure.\ref{fig:pdf_LAB}.

The second set of benchmarks simulates the actual business requirements of a telecommunication company in the 5G E2E network, denoted as APP. 
In APP, the distribution of each item is also randomly chosen in the set of uniform, truncated normal, fatigue-life, bimodal and gamma distribution, which is similar to the LAB case. 
However, in the 5G E2E network, each data package transmission is restricted to a time window with a fixed size, indicating that packets must be transmitted successfully within a set interval.
The difference between APP and LAB lies in the introduction of the retransmission mechanism \cite{nasrallah2018ultra}. 
This mechanism implies that each item in all classes has a probability of being retransmitted when a transmission failure occurs. 
In the APP, the first data package transmission attempt must be completed within $(0,10]$ ms with a success probability of 0.9. 
If the first attempt fails, a second attempt occurs with a delay of $(10,20]$ ms and a success probability of 0.9. 
Therefore, the delay of one transmission ranges from $(10,20]$ ms with a probability of $0.1\times 0.9 = 0.09$. 
The same applies to the third and fourth retransmissions.
The transmission fails if an item requires transmitting data more than four times, which limits the maximum number of retransmissions. 
Consequently, the time delay for an item to transmit a data package belongs to the intervals $(0,10],(10,20],(20,30],(30,40]$ with probabilities of $0.9,0.09,0.009,0.001$, respectively. 
The distribution of items involved in the APP is presented in Figure \ref{fig:pdf_APP}.
The cost generation formula for items is $c_{ij}=10/(\mu_{ij}+\sigma_{ij} )\times U(0.8,1.2)$. 
The purpose is to make the cost-effectiveness decrease as the delay decreases, thus simulating the marginal diminishing effect of cost-effectiveness in practical application.

For each benchmark in the LAB and APP, we design small-scale and large-scale problems with different parameter settings, as shown in Table \ref{benchmark: ss} and \ref{benchmark: ls} respectively. 
The scale of benchmarks, i.e., $m$ and $N$, is set similar to \cite{mkaouar2020solving}. 
The sample sizes $L$ are set 30 for small-scale benchmarks and 500 for large-scale benchmarks.
% The size of 30 is designed to simulate the situation where it is difficult to obtain a large number of samples in practical applications.
In Table \ref{benchmark: ss}, "\{11,14\}" of $T_{\max}^{LAB}$ means that $T_{\max}$ is set 11 or 14 as two different settings in the experiments of ss1. 

\begin{table}[htpb]
\centering
\caption{Small-scale benchmarks of DDCCMCKP}
\begin{tabular}{cccccc}
\toprule
\bfseries Benchmark & m & N & L & $T_{\max}^{LAB}$ & $T_{\max}^{APP}$\\
\midrule
{\bfseries $ss1$} &3 &5 &30 &\{11,14\} &\{21,27\}\\
{\bfseries $ss2$} &4 &5 &30 &\{18,26\} &\{47,49\}\\
{\bfseries $ss3$} &5 &5 &30 &\{11,20\} &\{27,38\}\\
{\bfseries $ss4$} &5 &10 &30 &\{10,16\} &\{16,27\}\\
\bottomrule
\end{tabular}
\label{benchmark: ss}
\end{table}

\begin{table}[htpb]
\centering
\caption{Large-scale benchmarks of DDCCMCKP}
\begin{tabular}{cccccc}
\toprule
\bfseries Benchmark & m & N & L & $T_{\max}^{LAB}$& $T_{\max}^{APP}$\\
\midrule
{\bfseries $ls1$} &10 &10 &500 &\{19,23\} &\{32,37\}\\
{\bfseries $ls2$} &10 &20 &500 &\{12,15\} &\{13,16\}\\
{\bfseries $ls3$} &20 &10 &500 &\{25,32\} &\{43,48\}\\
{\bfseries $ls4$} &30 &10 &500 &\{43,52\} &\{58,70\}\\
{\bfseries $ls5$} &40 &10 &500 &\{55,63\} &\{85,95\}\\
{\bfseries $ls6$} &50 &10 &500 &\{63,75\} &\{91,100\}\\
\bottomrule
\end{tabular}
\label{benchmark: ls}
\end{table}

\section{Computational Study}
\label{section:cs}
Experiments are carefully designed to evaluate the performance of DDALS in solving DDCCMCKP. 
We also conducted ablation studies, which involves disabling individual components of DDALS, to assess the contribution of the components integrated in DDALS.  
The source code of algorithms and the benchmark sets, are available at Github \footnote{Github: \href{https://github.com/eddylxf23/DDCCMCKP}{github.com/eddylxf23/DDCCMCKP}}. 
All experiments were run in python on an Intel(R) Xeon(R) CPU E5-2699A v4(2.40GHz,252G RAM) with Ubuntu 16.04 LST. 

\subsection{Algorithm Settings}

We set different parameter combinations of weighted parameter $\lambda$ and iteration times $t$ to evaluate the performance on small-scale and large-scale DDCCMCKP respectively. 
The performance is evaluated by the cost of final solution when the algorithm reaches the maximum of iteration, and the experiment was repeated for 10 times and take the average value.
% Since, $\lambda$ plays an important role in weights and utilities construction for items, $\lambda$ reflects the effect of the standard deviation of the sample data on the performance of DDALS. 
The value of $\lambda$ is tested from 1 to 10. 
As show in Fig.~\ref{fig:lambda}, there are several points worth noting. 
In the experiments of LAB-ss2, when $T_{\max} = 18$, the current solution of the algorithm remains unchanged. 
The results show that the lowest cost was obtained when $\lambda=0$. 
However, when we validate the output solution with real confidence level, we find that the solution was infeasible, which means that the algorithm could not find feasible solutions based on the samples without considering the standard deviation. 
When $T_{\max} = 26$, different lambdas do not affect the algorithm’s convergence to the lowest cost solution. 
In the experiments of LAB-ls2, when $T_{\max} = 25$, the algorithm converges very slowly when we set $\lambda=0$, while other settings can converge quickly. 
However, when $T_{\max} = 32$, different lambda values have little effect on the algorithm. 
That’s because the constraints are relatively loose, and feasible solutions account for a large proportion of the entire solution space, making the effect of standard deviation less significant. 
In general, it is necessary to set a non-zero value for $\lambda$, and our algorithm is robust to different $\lambda$, with $t=30$ being appropriate.

\begin{figure*}[!t]
\centering
\subfloat[LAB-ss2, $T_{\max} = 18$, exact]{\includegraphics[width=0.45\columnwidth]{chap6_exact_lambda_18.png}%
\label{fig:el18}}
\hfil
\subfloat[LAB-ss2, $T_{\max} = 26$, exact]{\includegraphics[width=0.45\columnwidth]{chap6_exact_lambda_26.png}%
\label{fig:el26}}
\hfil
\subfloat[LAB-ls3, $T_{\max} = 25$, MC]{\includegraphics[width=0.45\columnwidth]{chap6_MC_lambda_25.png}%
\label{fig:mcl25}}
\hfil
\subfloat[LAB-ls3, $T_{\max} = 32$, MC]{\includegraphics[width=0.45\columnwidth]{chap6_MC_lambda_32.png}%
\label{fig:mcl32}}
\caption{Performance analysis of different $\lambda$. As the algorithm iterates, lower cost can be achieved. The results show that the surrogate weight is necessary for DDALS while a non-zero $\lambda $ is necessary to the algorithm.}
\label{fig:lambda}
\end{figure*}

\subsection{Performance of Evaluation Methods}
Since each candidate solution can be seen as an unknown random distribution, concentration inequalities can serve as a reference to evaluate the quality of these candidate solutions. 
We utilize two kinds of concentration inequality as the lower bounds: Bernstein's Inequality and  $K$-order Hoeffding's Inequality, which have been widely used in the literature for evaluating the performance of optimization algorithms (Please refer to Appendix 1 for detailed formulas). 

\begin{table}[!t]
\caption{Evaluation for LAB-ss2}
\centering
\begin{tabular}{ccccccccccc}
\toprule
{\bfseries Confidence Level\%} & $S_1$ & $S_2$ & $S_3$ & $S_4$ & $S_5$\\
\midrule
{\bfseries Real} &68.24&89.41 &99.51 &99.98&100.00 \\
{\bfseries B's inequality} &23.22 &49.16 &75.63 &89.87 &95.85 \\
{\bfseries H's inequality} &22.83 &42.20 &62.32 &74.17 &85.12 \\
{\bfseries Exact} &81.64 &93.90 &99.53 &99.99 &100.00 \\
\bottomrule
\end{tabular}
\label{pef_eval_small}
\end{table}

We test exact method on LAB-ss2, and calculate the real confidence level for solution $S_1, S_2, S_3, S_4, S_5$ through the original generator, denoted as `Real'. 
`B's' and `H's' represent Bernstein's and Hoeffding's inequality, respectively. 
From Table \ref{pef_eval_small}, we found a substantial gap between the evaluation outcome and the real confidence level, indicating that the real distribution of solution of DDCCMCKP is far different from Gaussian distribution. 
Results show that the estimate of the exact method is closest to Real, and the estimation error becomes smaller as the real confidence level increases.


\begin{table}[!t]
\caption{Evaluation for APP-ls1}
\centering
\resizebox{\linewidth}{!}{
\begin{tabular}{ccccccccccc}
\toprule
{\bfseries Confidence Level\%} & $S_1$ & $S_2$ & $S_3$ & $S_4$ & $S_5$ \\
\midrule
{\bfseries Real} &79.18 &90.70 &94.48 &98.78 &98.81 \\
{\bfseries B's inequality} &9.56 &11.56 &19.82 &28.09 &48.87 \\
{\bfseries H's inequality} &12.58 &25.38 &38.30 &41.98 &53.96  \\
{\bfseries MC(10e2)} &81.1$\pm$ 3.78 &90.7$\pm$ 2.87&95.8$\pm$ 2.15&99.0$\pm$ 0.82&98.3$\pm$ 1.16 \\
{\bfseries MC(10e3)} &79.61$\pm$ 1.53&90.28$\pm$ 0.67&95.14$\pm$ 0.60&99.22$\pm$ 0.25&98.24$\pm$ 0.30 \\
{\bfseries MC(10e4)} &79.12$\pm$ 0.53&90.27$\pm$ 0.27&95.13$\pm$ 0.60&99.24$\pm$ 0.09&98.44$\pm$ 0.13  \\
{\bfseries MC(10e5)} &79.07$\pm$ 0.17&90.38$\pm$ 0.02&95.13$\pm$ 0.05&99.20$\pm$ 0.03&98.40$\pm$ 0.04 \\
{\bfseries MC(10e6)} &79.13$\pm$ 0.05&90.38$\pm$ 0.02&95.10$\pm$ 0.03&99.21$\pm$ 0.00&99.41$\pm$ 0.01  \\
\bottomrule
\end{tabular}
}
\label{pef_eval_large}
\end{table}

Next, we test Monte Carlo evaluation on APP-ls1 in Table \ref{pef_eval_large}. 
`MC(10e2)' represents that the evaluation is carried by 10e2 times of Monte Carlo estimation. 
The result is the average of 10 times' results. 
`81.1 $\pm$ 3.78' represents that the mean of estimated confidence level is 81.1\% and the standard deviation is 3.78\%. 
The result shows that the estimation error of $10e2$ times has the largest fluctuations and $10e5$ tends to be close to the real confidence value, but still not enough to meet the requirement for larger scale of problem and higher confidence level. 
Considering the trade-off between accuracy and speed, we set $10e6$ as the Monte Carlo simulation times.

\subsection{Results on Small-Scale and Large-Scale DDCCMCKP}
To evaluate the effectiveness of DDALS, we compare DDALS with several baseline methods on various benchmark sets from the real-world application. 
The baseline methods include Greedy, GA and EDA. 
Greedy typically works well on problems where the locally optimal choice leads to an overall optimal solution.
GA performs well in solving different optimization problems \cite{9075201}.
While EDA achieves great success in both combinatorial and continuous optimization \cite{6899662}. 
% In the experiments, the design of Greedy refers to \cite{kellerer2004multiple}.
% For GA, the population size is 10, the probability of cross-over is 0.1 and the size of elite population is 6. 
% As for EDA, the population size is 10 and the size of elite population is 6. 

% Each experiment will be run 10 times, and the results will be recorded as the mean $\pm$ standard deviation.
% We use evaluation times (ET) as the metric to evaluate the algorithm’s running speed, since most of the time consumed by the algorithm comes from solution evaluation. 
For fairness, we first run DDALS then record the value of ET. 
Next, we let the Greedy, GA, and EDA keep running until convergence or the ET of which exceeds that of DDALS, then output the current solution as the result. 
The population size is set to 10 and the number of dominant individuals is 6 for both GA and EDA. 
For GA, the crossover probability of parents is set to 0.1 and the probability of variation is $1/m$. 
The real confidence level is calculated by using the original generators to generate $10e7$ data. 
Note that a solution is regarded as feasible when its real confidence level is larger than $0.985$, which means that we allow for an error margin of $0.5\%$.


% C: Cost, ET: Evaluation times, CL: confidence level, RCL: real CL, FSR: feasible solution ratio(in 10 trials)
\begin{table*}[!t]
    \caption{Performance on benchmarks of LAB. The highest average real confidence level and the lowest cost are highlighted through bolding and being underlined.}
    \centering
    \resizebox{\linewidth}{!}{
    \begin{tabular}{ccccccccccc}
    \toprule
    \multirowcell{2}{\bfseries Benchmark} & \multicolumn{5}{c}{\bfseries small $T_{\max}$} & \multicolumn{5}{c}{\bfseries large $T_{\max}$} \\
    \cmidrule(r){2-6} \cmidrule(r){7-11}
    & {\bfseries C} & {\bfseries ET} & {\bfseries ECL} & {\bfseries RCL} & {\bfseries FSR} & {\bfseries C} & {\bfseries ET} & {\bfseries ECL} & {\bfseries RCL} & {\bfseries FSR}\\
     
    \midrule
    {\bfseries ss3-Greedy} &24.4$\pm$0.0&2.0$\pm$0.0 & 100$\pm$0&99.9$\pm$0.0 &\underline{\bfseries 1} &21.8$\pm$0.0&3.0$\pm$0.0& 100$\pm$0&99.9$\pm$0.0 &\underline{\bfseries 1} \\
    {\bfseries ss3-GA} &22.6$\pm$1.5&303.4$\pm$2.8 & 99.9$\pm$0.0 & 99.7$\pm$0.2 &\underline{\bfseries 1} &12.9$\pm$0.7 & 262.4$\pm$1.6&99.9$\pm$0.0&99.9$\pm$0.1 &\underline{\bfseries 1} \\
    {\bfseries ss3-EDA} &24.1$\pm$2.3&304.1$\pm$2.3 & 99.9$\pm$0.1&99.6$\pm$0.6&0.9&14.0$\pm$2.0&264.5$\pm$2.9& 99.9$\pm$0.0&99.9$\pm$0.1 &\underline{\bfseries 1} \\
    {\bfseries ss3-DDALS(O)} &\underline{\bfseries 21.4$\pm$0.0} &291.9$\pm$6.2 & 99.9$\pm$0&99.5$\pm$0.0 &\underline{\bfseries 1} &\underline{\bfseries 12.6$\pm$0.0} &254.4$\pm$4.7& 100$\pm$0 &99.9$\pm$0  &\underline{\bfseries 1} \\
    {\bfseries ss3-DDALS(V1)} &\underline{\bfseries 21.4$\pm$0.0} &291.9$\pm$6.2 & 99.9$\pm$0&99.5$\pm$0.0 &\underline{\bfseries 1} &\underline{\bfseries 12.6$\pm$0.0}&254.4$\pm$4.7& 100$\pm$0&99.9$\pm$0  &\underline{\bfseries 1} \\ 
    {\bfseries ss3-DDALS(V2)} &\underline{\bfseries 21.4$\pm$0.0} &291.9$\pm$6.2 & 99.9$\pm$0&99.5$\pm$0.0 &\underline{\bfseries 1} &\underline{\bfseries 12.6$\pm$0.0}&254.4$\pm$4.7& 100$\pm$0&99.9$\pm$0  &\underline{\bfseries 1} \\
    {\bfseries ss3-DDALS(V3)} &23.9$\pm$1.7&279 & 99.8$\pm$0.2&99.1$\pm$0.8&0.8&15.7$\pm$1.9&260& 99.9$\pm$0.3&99.6$\pm$0.9  &0.9 \\
    
    \midrule
    {\bfseries ss4-Greedy} &30.1$\pm$0.0&2.0$\pm$0.0 & 99.9$\pm$0.0&99.9$\pm$0.0 &\underline{\bfseries 1} &27.5$\pm$0.0&3.0$\pm$0.0& 99.7$\pm$0.0&99.4$\pm$0.0  &\underline{\bfseries 1} \\
    {\bfseries ss4-GA} &23.8$\pm$5.5&768.9$\pm$8.5 & 99.9$\pm$0.3&99.7$\pm$0.4 &\underline{\bfseries 1} &17.2$\pm$3.0&562.6$\pm$5.7& 99.7$\pm$0.3&99.6$\pm$0.3  &\underline{\bfseries 1} \\
    {\bfseries ss4-EDA} &\underline{\bfseries 17.4$\pm$0.0} &763.4$\pm$2.5 & 99.9$\pm$0.0&99.9$\pm$0.0 &\underline{\bfseries 1} &13.0$\pm$1.0&553.4$\pm$2.7& 99.8$\pm$0.2&98.9$\pm$0.4  &\underline{\bfseries 1} \\
    {\bfseries ss4-DDALS(O)} &\underline{\bfseries 17.4$\pm$0.0} &759.9$\pm$25.8 & 99.9$\pm$0.0&99.9$\pm$0.0 &\underline{\bfseries 1} &\underline{\bfseries 11.2$\pm$0.0} &525.3$\pm$39.1& 99.2$\pm$0.0&99.0$\pm$0.0  &\underline{\bfseries 1} \\
    {\bfseries ss4-DDALS(V1)} &\underline{\bfseries 17.4$\pm$0.0} &759.9$\pm$25.8 & 99.9$\pm$0.0&99.9$\pm$0.0 &\underline{\bfseries 1} &12.1$\pm$0.0&525.3$\pm$39.1& 99.5$\pm$0.0&98.9$\pm$0.0  &\underline{\bfseries 1} \\
    {\bfseries ss4-DDALS(V2)} &\underline{\bfseries 17.4$\pm$0.0} &759.9$\pm$25.8 & 99.9$\pm$0.0&99.9$\pm$0.0 &\underline{\bfseries 1} &14.0$\pm$1.2&525.3$\pm$39.1& 99.9$\pm$0.2&99.0$\pm$0.3  &\underline{\bfseries 1} \\
    {\bfseries ss4-DDALS(V3)} &24.6$\pm$4.0&735 & 99.8$\pm$0.2&99.5$\pm$0.5 &\underline{\bfseries 1} &13.8$\pm$1.3&561& 99.8$\pm$0.3&99.5$\pm$0.5  &\underline{\bfseries 1} \\
    
    \midrule
    {\bfseries ls4-Greedy} &169.7$\pm$0.0&5$\pm$0.0 & 99.8$\pm$0.0&99.8$\pm$0.0 &\underline{\bfseries 1} &162.5$\pm$0.0&6.0$\pm$0.0& 99.9$\pm$0.0&99.9$\pm$0.0  &\underline{\bfseries 1} \\
    {\bfseries ls4-GA} &123.1$\pm$8.8&11307.2$\pm$8.0 & 99.1$\pm$0.1&99.1$\pm$0.2 &\underline{\bfseries 1} &104.1$\pm$5.1&8409.1$\pm$5.6& 99.1$\pm$0.2&99.2$\pm$0.2  &\underline{\bfseries 1} \\
    {\bfseries ls4-EDA} &126.2$\pm$2.3&11304.8$\pm$2.1 & 99.2$\pm$0.2&99.4$\pm$0.2 &\underline{\bfseries 1} &114.8$\pm$4.5& 8403.4$\pm$2.8&99.6$\pm$0.2&99.6$\pm$0.2  &\underline{\bfseries 1} \\
    {\bfseries ls4-DDALS(O)} &\underline{\bfseries 106.6$\pm$0.8} &11289.9$\pm$92.0 & 99.1$\pm$0.1&99.2$\pm$0.1 &\underline{\bfseries 1} &\underline{\bfseries 92.8$\pm$0.7} &8376.6$\pm$378.8& 99.1$\pm$0.1&99.2$\pm$0.1  &\underline{\bfseries 1} \\
    {\bfseries ls4-DDALS(V1)} &109.5$\pm$0.4&11289.9$\pm$92.0 & 99.7$\pm$0.1&99.7$\pm$0.1 &\underline{\bfseries 1} &97.9$\pm$2.2&8376.6$\pm$378.8& 99.6$\pm$0.1&99.7$\pm$0.1  &\underline{\bfseries 1} \\
    {\bfseries ls4-DDALS(V2)} &108.9$\pm$1.3&11289.9$\pm$92.0 & 99.6$\pm$0.2&99.7$\pm$0.2 &\underline{\bfseries 1} &93.8$\pm$1.4&8376.6$\pm$378.8& 99.3$\pm$0.2&99.4$\pm$0.2  &\underline{\bfseries 1} \\
    {\bfseries ls4-DDALS(V3)} &108.7$\pm$1.1&11171 & 99.2$\pm$0.2&99.3$\pm$0.2 &\underline{\bfseries 1} &96.5$\pm$1.7&8177& 99.1$\pm$0.1& 99.2$\pm$0.1  &\underline{\bfseries 1} \\
    
    \bottomrule
    \end{tabular}
    }
\label{result:LAB}
\end{table*}

The results on LAB benchmarks are presented in Table \ref{result:LAB}. 
`ss3-Greedy' in the left of table indicates the result of Greedy tested on benchmark ss3. 
`DDALS(O)', `DDALS(V1)', `DDALS(V2)' and `DDALS(V3)' represents the original, first, second and third version of DDALS. 
`small $T_{\max}$' and `large $T_{\max}$' are corresponding to small and large $T_{\max}$ in Table \ref{benchmark: ss} and Table \ref{benchmark: ls}.
Five indicators are utilized to measure the performance. 
The cost of solution is denoted as `C'. 
The evaluation times `ET' is utilized to compare the convergence speed of algorithms. 
`ECL' indicates estimated confidence level, while `RCL' is calculated by real generator. 
`FSR'  represents the ratio of actual feasible solution in the output results of the 10 experiments. 
All results are based on statistical analysis from 10 repetitions of the experiment, denoted as `mean $\pm$ standard deviation'.
Due to space limitation, parts of results are included in the main text, and the complete results are available in Appendix 5.

\begin{table*}[!t]
    \caption{Performance on benchmarks of APP. Due to the introduction of retransmission mechanisms, there is a significant difference in the results between LAB and APP.}
    \centering
    \resizebox{\linewidth}{!}{
        \begin{tabular}{ccccccccccc}
        \toprule
        \multirowcell{2}{\bfseries Benchmark} & \multicolumn{5}{c}{\bfseries small $T_{\max}$} & \multicolumn{5}{c}{\bfseries large $T_{\max}$} \\
        \cmidrule(r){2-6} \cmidrule(r){7-11}
        & {\bfseries C} & {\bfseries ET} & {\bfseries ECL} & {\bfseries RCL} & {\bfseries FSR} & {\bfseries C} & {\bfseries ET} & {\bfseries ECL} & {\bfseries RCL} & {\bfseries FSR}\\
        
        \midrule
        {\bfseries ss2-Greedy} & 6.4$\pm$0.0 & 3$\pm$0.0 & 99.7$\pm$0.0 &98.6$\pm$0.0 &\underline{\bfseries 1} &6.4$\pm$0.0&3.0$\pm$0.0& 99.8$\pm$0.0&99.0$\pm$0.0  &\underline{\bfseries 1} \\
        {\bfseries ss2-GA} &5.3$\pm$0.0&301.8$\pm$2.2 & 99.2$\pm$0.0&97.3$\pm$0.4&0&5.3$\pm$0.0&272.6$\pm$2.8& 99.5$\pm$0.1&98.1$\pm$0.2  &0 \\
        {\bfseries ss2-EDA} &5.4$\pm$0.2&303.1$\pm$3.0 & 99.3$\pm$0.2&97.4$\pm$0.7&0.1&5.4$\pm$0.2&274.9$\pm$3.0& 99.6$\pm$0.2&98.3$\pm$0.4  &0.3 \\
        {\bfseries ss2-DDALS(O)} &\underline{\bfseries 5.2$\pm$0.0} &307$\pm$7.6 & 99.2$\pm$0.0&97.3$\pm$0&0 &\underline{\bfseries 5.2$\pm$0.0} &291$\pm$10.6& 99.6$\pm$0.0&98.1$\pm$0.0  &0 \\
        {\bfseries ss2-DDALS(V1)} &5.5$\pm$0.1&307$\pm$7.6 & 99.8$\pm$0.0&97.5$\pm$0.3&0 &\underline{\bfseries 5.2$\pm$0} &291$\pm$10.1& 99.6$\pm$0.0&98.1$\pm$0.0  &0 \\
        {\bfseries ss2-DDALS(V2)} &5.5$\pm$0.1&307$\pm$ 7.6 & 99.8$\pm$0.0&97.5$\pm$0.3&0&5.4$\pm$0.1&291$\pm$10.1& 99.9$\pm$0.1&98.2$\pm$0.1  &0 \\
        {\bfseries ss2-DDALS(V3)} &5.4$\pm$0.1&295 & 99.4$\pm$0.3&97.7$\pm$0.4&0&5.3$\pm$0.1&272.6$\pm$2.8& 99.5$\pm$0.1&98.1$\pm$0.2  &0 \\
        
        \midrule
        {\bfseries ss4-Greedy} &34.9$\pm$0.0&3.0$\pm$0.0 & 100.0$\pm$0.0&100.0$\pm$0.0 &\underline{\bfseries 1} &27.2$\pm$0.0&4.0$\pm$0.0& 100.0$\pm$0.0&99.0$\pm$0.0  &\underline{\bfseries 1} \\
        {\bfseries ss4-GA} &27.3$\pm$1.1&2010$\pm$8.4 & 99.8$\pm$0.3&90.4$\pm$0.5&0&21.5$\pm$2.9&1614.8$\pm$9.8& 99.6$\pm$0.3&93.2$\pm$4.8  &0 \\
        {\bfseries ss4-EDA} &\underline{\bfseries 26.1$\pm$0.0} &2005.7$\pm$2.3 & 100$\pm$0&90.1$\pm$0&0 &\underline{\bfseries 16.4$\pm$0.0} &1602.7$\pm$2.1& 99.9$\pm$0.0&80.3$\pm$0.0  &0 \\
        {\bfseries ss4-DDALS(O)} &\underline{\bfseries 26.1$\pm$0.0} &1988.2$\pm$30.4 & 1$\pm$0&90.1$\pm$0&0&16.7$\pm$0.9&1579.4$\pm$76.8 & 99.8$\pm$0.2&81.8$\pm$4.4&0 \\
        {\bfseries ss4-DDALS(V1)} &\underline{\bfseries 26.1$\pm$0.0} &1988.2$\pm$30.4 & 1$\pm$0&90.1$\pm$0&0&16.7$\pm$1.0&1579.4$\pm$ 76.8& 99.9$\pm$0.0&81.5$\pm$3.6  &0 \\
        {\bfseries ss4-DDALS(V2)} &\underline{\bfseries 26.1$\pm$0.0} &1988.2$\pm$30.4 & 1$\pm$0&90.1$\pm$0&0&18.0$\pm$1.9&1579.4$\pm$76.8& 99.9$\pm$0.1&86.1$\pm$7.0  &0 \\
        {\bfseries ss4-DDALS(V3)} &29.1$\pm$2.4&2022 & 99.9$\pm$0.2&91.3$\pm$2.9&0&19.9$\pm$1.7&1589& 99.7$\pm$0.3&92.0$\pm$5.3  &0 \\
        
        \midrule
        {\bfseries ls2-Greedy} &85.7$\pm$0.0&3.0$\pm$0.0 & 99.9$\pm$0.0&99.9$\pm$0.0 &\underline{\bfseries 1} &85.7$\pm$0.0&3.0$\pm$0.0& 100.0$\pm$0.0&100.0$\pm$0.0  &\underline{\bfseries 1} \\
        {\bfseries ls2-GA} &76.9$\pm$ 2.7&22617.4$\pm$10.9 & 99.0$\pm$0.1&98.9$\pm$0.2 &\underline{\bfseries 1} &61.2$\pm$2.1&21327.1$\pm$20.3& 99.1$\pm$0.0&99.1$\pm$0.1  &\underline{\bfseries 1} \\
        {\bfseries ls2-EDA} &75.9$\pm$0.0&22604.3$\pm$3.3 & 99.0$\pm$0.0&98.9$\pm$0.0 &\underline{\bfseries 1} &59.2$\pm$0.3&21306.5$\pm$2.5& 99.2$\pm$0.0&99.2$\pm$0.1  &\underline{\bfseries 1} \\
        {\bfseries ls2-DDALS(O)} &\underline{\bfseries 71.9$\pm$0.3} &22591.3$\pm$63.6 & 99.0$\pm$0.0&99.0$\pm$0.0 &\underline{\bfseries 1} &\underline{\bfseries 57.2$\pm$0.0} &21233.5$\pm$51.9& 99.3$\pm$0.1&99.3$\pm$0.1  &\underline{\bfseries 1} \\
        {\bfseries ls2-DDALS(V1)} &76.6$\pm$1.1&22591.3$\pm$63.6 & 99.6$\pm$0.1&99.5$\pm$0.1 &\underline{\bfseries 1} &58.5$\pm$0.8&21233.5$\pm$51.9& 99.6$\pm$0.1&99.6$\pm$0.1  &\underline{\bfseries 1} \\
        {\bfseries ls2-DDALS(V2)} &72.9$\pm$0.6&22591.3$\pm$63.6 & 99.3$\pm$0.2&99.2$\pm$0.2 &\underline{\bfseries 1} &57.5$\pm$0.2&21233.5$\pm$51.9& 99.4$\pm$0.1&99.5$\pm$0.1  &\underline{\bfseries 1} \\
        {\bfseries ls2-DDALS(V3)} &73.1$\pm$0.7&22564 & 99.2$\pm$0.1&99.1$\pm$0.1 &\underline{\bfseries 1} &57.9$\pm$0.4&21245& 99.2$\pm$0.1&99.2$\pm$0.2  &\underline{\bfseries 1} \\
        
        \midrule
        {\bfseries ls4-Greedy} &207.3$\pm$0.0&13.0$\pm$0.0 & 99.0$\pm$0.0&98.9$\pm$0.0 &\underline{\bfseries 1} &195.3$\pm$0.0&17.0$\pm$0.0& 99.4$\pm$0.0&99.4$\pm$0.0  &\underline{\bfseries 1} \\
        {\bfseries ls4-GA} &208.5$\pm$8.3&48014.1$\pm$12.1 & 99.0$\pm$0.0&98.6$\pm$0.3&0.7&175.2$\pm$10.3&45014.6$\pm$10.2& 99.0$\pm$0.0&98.5$\pm$0.2  &0.6 \\
        {\bfseries ls4-EDA} &205.3$\pm$2.4&48004.5$\pm$3.3 & 99.0$\pm$0.0&99.0$\pm$0.0 &\underline{\bfseries 1} &163.8$\pm$1.3&45005.5$\pm$2.6& 99.1$\pm$0.0&98.8$\pm$0.0  &\underline{\bfseries 1} \\
        {\bfseries ls4-DDALS(O)} &\underline{\bfseries 194.9$\pm$0.8} &47962.8$\pm$67.7 & 99.0$\pm$0.0&98.9$\pm$0.0 &\underline{\bfseries 1} &\underline{\bfseries 163.3$\pm$3.6} &44822.1$\pm$583.5& 99.0$\pm$0.0&98.4$\pm$0.1  &0.3 \\
        {\bfseries ls4-DDALS(V1)} &203.2$\pm$2.0&47962.8$\pm$67.7 & 99.1$\pm$0.0&99.0$\pm$0.1 &\underline{\bfseries 1} &172.8$\pm$4.5&44822.1$\pm$583.5& 99.3$\pm$0.1&98.9$\pm$0.3  &0.8 \\
        {\bfseries ls4-DDALS(V2)} &196.5$\pm$1.3&47962.8$\pm$67.7 & 99.0$\pm$0.0&99.0$\pm$0.0 &\underline{\bfseries 1} &165.3$\pm$4.3&44822.1$\pm$583.5& 99.1$\pm$0.1&98.5$\pm$0.2  &0.3 \\
        {\bfseries ls4-DDALS(V3)} &196.9$\pm$2.1&48010 & 99.0$\pm$0.0&98.9$\pm$0.1 &\underline{\bfseries 1} &168.2$\pm$ 2.3&44174& 99.0$\pm$0.0&98.9$\pm$0.1  &\underline{\bfseries 1} \\
        
        \bottomrule
        \end{tabular}
    }
    \label{result:APP}
\end{table*}

In the case of LAB scenarios, the results of ss3, ss4 and ls4 are shown in Table \ref{result:LAB}. 
We can intuitively see that the DDALS algorithm outputs the lowest cost solution across all test benchmarks, while Greedy needs only less than 10 times of evaluation but results in the highest cost. 
The feasibility ratio of all algorithms is close to 1, indicating that our solution evaluation algorithm only needs a very few samples to obtain an evaluation that is close to the true confidence level for small-scale scenarios and 500 samples for large-scale scenarios.
From another perspective, for other variations, the original version of DDALS can obtain solutions with lower cost than other variations while guaranteeing a feasible ratio of 1, indicating that the solution evaluation is accurate. 

There are significant changes in the scenario of APP, shown in Table \ref{result:LAB}. 
In small-scale benchmark tests, only the Greedy achieves a feasibility ratio of 1, while the feasibility ratio of other algorithms is almost zero. 
Although DDALS can obtain lower cost solutions, all of them are found to be infeasible. 
Such results are practically useless. 
The main reason is due to the retransmission mechanism in APP, which introduces discontinuities in the item distribution that pose a significant challenge to the solution evaluation algorithm. 
In a random sampling case, small samples may easily miss information regarding exceptional jumps, leading to performance decline of solution evaluation.
Greedy can achieve higher feasibility ratios, but it sacrifices the optimality of the solutions. 
For the results of large-scale problems, although the difficulty of search increased with the increase of scale,  DDALS can still find feasible solutions achieving the lowest cost.
Note that in the result of ls4 with large $T_{\max}$, the feasibility ratio of DDALS(O) is 0.3 while that of DDALS(V3) is 1. 
Given that the average quality of the solutions generated by DDALS(V3) decreases by only 3\% relative to the original DDALS, the feasibility ratio increased by 233\%.
It shows the effectiveness and necessity of filtering the solutions in complex scenarios. 
Moreover, increasing the sample size $L$ can improve the evaluation performance, which can help capture more information and improve the accuracy of solution evaluation.

\subsection{Effectiveness of Components in DDALS}
We conduct an ablation study to further assess the effectiveness of four components of DDALS, including CP, LSS, Degrade and FSS. 
To independently assess their effectiveness, four different variants of DDALS are proposed and tested on benchmarks of LAB:

\begin{enumerate}
\item r-CP DDALS: It uses random strategy to select the item in each class instead of utilizing CP. 
Note that if the initial randomly generated is infeasible solution, query the current solution for each class of item and re-generate randomly in each class of items less than the weight of the current item.
\item no-LSS DDALS: The LSS procedure in the iteration process is removed from DDALS.
\item no-Degrade DDALS: The degrade procedure in the iteration process is removed. 
\item no-FSS DDALS: The FSS procedure is removed. 
\end{enumerate}

For all the above variants, we compute the performance degradation ratio(PDR) by computing $(C_2-C_1)/C1$ for each benchmark, where $C_1$ and $C_2$ are the average cost of solution obtained by the variant and original across 10 independent runs. 
The component being examined is useful for a specific LAB benchmark if the result shows that PDR $>0$. 
Table \ref{variant of LS} presents the average PDR averaged on each benchmark of LAB. 
Based on the results, it can be concluded that none of the examined components has a negative effect on the performance of DDALS.
In most cases, these components make LS perform better. 
Constructing utility for items is helpful for DDALS to obtain a good initial solution on average. 
LSS provides some kind of exploration, but a negative value on ls3 implies that due to the large size of the sample, the algorithm may still not converge after 30 iterations of running. 
The operation of degradation has the most significant contribution to the performance of LS, since it greatly improves the problem of the algorithm, as it helps the algorithm to escape local optima. 
A deeper search is necessary, while FSS has a slightly higher PDR than LSS. 
However, with the increase of the benchmark's scale, the contribution of FSS seems to decrease, implying that more efficient and exploratory methods are needed to cope with larger-scale problems.

\begin{table}[!t]
\centering
\caption{Average PDR on each benchmark of LAB. The highest PDR is highlighted through bolding and being underlined. }
\begin{tabular}{ccccc}
\toprule
\bfseries Benchmark & r-CP & no-LSS & no-Degrade & no-FSS  \\
\midrule
{\bfseries $ss1$} &0.00\% &0.00\% &0.00\% &0.00\% \\
{\bfseries $ss2$} &0.00\% &0.00\% &\underline{\bfseries 4.81\%} &0.00\% \\
{\bfseries $ls1$} &0.16\% &0.16\% &0.16\% &\underline{\bfseries 0.73\%} \\
{\bfseries $ls2$} &0.10\% &\underline{\bfseries 1.07\%} &0.63\% &0.74\% \\
{\bfseries $ls3$} &-0.21\% &0.63\% &\underline{\bfseries 0.70\%} &0.54\% \\
{\bfseries $ls4$} &\underline{\bfseries 1.07\%} &-0.69\% &0.09\% &\underline{\bfseries 0.22\%} \\
\midrule
{\bfseries Avg.PDR} &0.19\% &0.19\% &\underline{\bfseries 1.06\%} &0.69\% \\
\bottomrule
\end{tabular}
\label{variant of LS}
\end{table}

\subsection{Acceleration Performance of Accelerated Monte Carlo}
The acceleration effect of Accelerated Monte Carlo (AMC) evaluation is validated. 
We randomly selected 10,000 solutions for evaluation, using regular Monte Carlo (MC) and AMC respectively. 
After recording the running time required for each 1000 evaluations, the result is shown in Figure \ref{accelerate}. 

The results show that the time consumed increases essentially linearly as the number of evaluation times increases. 
For different $T_{\max}$ settings, MC takes essentially the same amount of time since the number of simulations needed does not vary with $T_{\max}$, while AMC can obtain better acceleration compared to MC when $T_{\max}$ is small. 
For example, compared with the running time $526.601s$ of MC, AMC takes $458.449s$ to finish the evaluation when $T_{\max}=11$, while the number of fast-filtered solution is 10000. 
When $T_{\max}=23$, it takes AMC about $471.395s$, although the number of pre-fast filtered is 9718. 
As $T_{\max}$ increases, the difficulty of the chance constraint decreases, resulting in the increase of the percentage of feasible solutions in the whole solution space. 
Thus, the number of solutions that can be fast pre-filtered by sequences decreases, which leads to a decrease in the acceleration effect. 
The accelerated mechanism may have a negative impact on the overall efficiency of the evaluation as $T_{\max}$ becomes large. 
When $T_{\max}=27$, AMC and MC take almost the same time, while when $T_{\max}=35$, AMC needs to take $807.435s$ to finish rather than about $500s$ of MC. 
In conclusion, when applying the accelerated mechanism, a pre-assessment of whether the constraints are tight enough is necessary to determine whether the constraints are tight enough.

\begin{figure}[!t]
\centering
\includegraphics[width=3in]{chap6_Accelarate_performance.png}
\caption{Acceleration performance of accelerated Monte Carlo on LAB-ls1. The mechanism has a positive effect when the constraints are tight, but has a negative impact on the overall efficiency as $T_{\max}$ becomes large.} \label{accelerate}
\end{figure}

\subsection{Performance of Other Works}
\label{subsection:jy}
We also attempted to adjust two similar algorithms, Ji's \cite{ji2021data} and Yang's \cite{yang2018algorithm}, to solve our problem (see Appendix 4 for modification details). Due to the incapability of dynamic programming and mixed integer programming for solving large-scale DDCCMCKP, we conducted experiments on small-scale benchmarks to evaluate their performance. 
Results are shown in Table \ref{Result MILP and iter-MCKP}, and we observed that the majority of the results were either "infeasible" or "fail". 
"fail" means that the algorithms fail to output any solution, while "infeasible" means that the optimal solution is infeasible. 
The value of "Feasibility" column is the total cost of solution. 
For Yang's method,  the poor performance was mainly due to the significant difference between the actual distribution and the Gaussian distribution, and the insufficient data to characterize an unknown distribution using sample mean and variance, leading to incorrect iteration.
In Ji's work, a uniform distribution was used and tested. The performance of their method is poor since it could not handle the long-tailed distribution of the candidate solutions in our problem and led to suboptimal performance.

\begin{table}[!t]
\centering
\caption{Result of Ji's \cite{ji2021data} and Yang's \cite{yang2018algorithm}}
\resizebox{\linewidth}{!}{
\begin{tabular}{ccccc}
\toprule
\multirowcell{2}{\bfseries Benchmark} & \multicolumn{2}{c}{\bfseries Ji} & \multicolumn{2}{c}{\bfseries Yang} \\
\cmidrule(r){2-3} \cmidrule(r){4-5}
& {\bfseries Time(s)} & {\bfseries Feasiblility} & {\bfseries Time(s)} & {\bfseries Feasiblility}\\
\midrule
{\bfseries LAB-$ss1$} &- &{\bfseries fail} &1.0 &14.5\\
{\bfseries LAB-$ss2$} &0.4 &{\bfseries infeasible} &- &{\bfseries fail}\\
{\bfseries LAB-$ss3$} &0.8 &{\bfseries infeasible} &1.4 &21.5\\
{\bfseries LAB-$ss4$} &16 &{\bfseries infeasible} &43.2 &17.4\\
{\bfseries APP-$ss1$} &- &{\bfseries fail} &0.1 &{\bfseries infeasible}\\
{\bfseries APP-$ss2$} &0.6 &{\bfseries infeasible} &- &{\bfseries fail}\\
{\bfseries APP-$ss3$} &- &{\bfseries fail} &0.2 &{\bfseries infeasible}\\
{\bfseries APP-$ss4$} &- &{\bfseries fail} &13.5 &{\bfseries infeasible}\\
\bottomrule
\label{Result MILP and iter-MCKP}
\end{tabular}
}
\end{table}

\section{Conclusion}
\label{section:con}
In this work, derived from practical applications, we study a new variant of MCKP called DDCCMCKP. 
To address this specific problem, a novel DDALS algorithm has been introduced.
Based on local search framework, DDALS enables data-driven solution evaluation for the problem,  and incorporates a surrogate weight to ensure the optimality of solutions. 
Moreover, a solution filtering mechanism is introduced to improve the feasibility of the output solutions. 
Two sets of benchmarks for DDCCMCKP have been proposed, one of which simulates the real business requirement of 5G E2E network scenario for a communication company. 
The experimental results demonstrate that DDALS is able to solve DDCCMCKP effectively and outperforms other baseline methods.
Although the APP benchmark imposes higher demands on problem-solving abilitiy, the performance of DDALS is not satisfactory in small-scale problems. 
However, DDALS can solve DDCCMCKP well in most of the scenarios and maintain better performance than other baselines.

Overall, our proposed algorithm holds great promise for addressing this new variant of MCKP in the data-driven scenario. 
DDALS can serve as a baseline for future research. The code and benchmark sets will be open-sourced for further promoting research on DDCCMCKP.
The algorithm framework can be further improved in two aspects. 
A new surrogate weight design may improve the efficiency and solution's quality instead of the current design.
From the other perspective, more effective filtering strategies that can balance the feasibility and optimality of the solutions could be explored in the future.


\appendices
\section{Details about Concentration Inequalities}
\subsection{Bernstein's Inequality}
Assume that there are $m$ classes at the E2E, each denoted as variable $X_j,j=1,\ldots,m$. 
$X_1,X_2,\ldots,X_m$ are independent and satisfy for given some $\{(a_j,b_j),\forall{j}\in\{1,\ldots,m\} \}$ and constant $C$ there are $P(X_j\in [a_j,b_j])=1, \forall{j}\in\{1,\ldots,m\}$ and $b_j-a_j=c_j\leq C$. 
Denote $S=X_1+X_2+\ldots+X_m$ as the sum of random variables $X_j$ and $V=\sum_{i=1}^{m}Var(X_j)$ is the sum of variances. 
Then we have Bernstein's inequality,

\begin{equation}\label{eqa:Bernstein}
    P(S\leq E(S)+\alpha)\geq 1-exp(-\frac{\alpha^2/2}{V+C\alpha /3})
\end{equation}

\subsection{K-order Hoeffding's Inequality}
Similar to Bernstein's inequality, we assume that $X_1,X_2,\ldots,X_m$ are independent and satisfy for given some $a_j,b_j$, there is $P(X_j\in [a_j,b_j])=1, \forall{j}\in\{1,\ldots,m\}$. 
Denote new variable $Y_j=X_j-a_j$, which satisfies $P(Y_j\in [0,c_j])=1,\forall{j}\in\{1,\ldots,m\}$. 
Denote $S=Y_1+Y_2+\ldots+Y_m$ and we have $\mathbb{E}(Y_j^k)=\mu _j^k$. $K$ is the order of origin moment. 
Thus, we have following $K$-order Hoeffding's inequality:

\begin{equation}\label{eqa:Hoeffding}
\begin{aligned}
    & P(S\leq E(S)+\alpha)\geq \\
    & 1-exp(-\frac{2\alpha^2}{\sum_{j=1}^m b_j^2 C_K(\frac{4\alpha b_j}{D},b_j,\mu_j^1,\ldots,\mu_j^K)})
\end{aligned}
\end{equation}
where $D=\sum_{j=1}^m(\mu_j^2/\mu_j^1)^2$ and $C_K(\frac{4\alpha b_j}{D},b_j,\mu_j^1,\ldots,\mu_j^K)$ is a value calculated by all the orders of origin moments of $Y_j$.
However, the concentration inequality does not make any assumptions about the form of the distribution of the variables; hence, the lower bound provided by the inequality may be imprecise. Additionally, as the inequalities are solely based on data-driven background, they may only serve as rough estimates for the upper and lower bounds of the variables and this may introduce extra errors into the confidence evaluation.

Since only data of items are available, if we manage to calculate the ratio of the combination of data meeting the upper bound $W$, to the total number of data combinations, then we can know whether the solution satisfies the chance constraint. 
For uncertain solution $S=[s_1,s_2,\ldots,s_m]$, we count the number of $\sum_{i=1}^{m}d_{i,j_i,l}\leq W$ and check if it exceeds $P_0L^m$. 
Since $P_0$ is usually a number very close to 1 under high confidence level requirement, it is more practical to  compare the number of $\sum_{i=1}^{m}d_{s_il}\geq W$ with $(1-P_0)L^m$. 
If the number of $\sum_{i=1}^{m}d_{s_il}\geq W$ is greater than $(1-P_0)L^m$, then we can determine quickly that it is an infeasible solution.

\section{Proof}

\begin{theorem}
The popped sum in the $k^{th}$ iteration is exactly the $k^{th}$ largest sum.
\end{theorem}

\begin{proof}
Suppose the $i^{th}$ largest sum is $S_r = d_{s_1l_1}+d_{s_2l_1}+\ldots+d_{s_ml_m}$.
\begin{itemize}
\item For any pair in the combination or queue $Q$ that has already been output, the combinations with indices less than or equal to it must have been popped  or exist in $Q$.
\item   Construct a chain $S$, then we have the following chain, denoted as $S_0\leq S_1\leq ,\ldots,\leq S_l$.
\begin{equation*}
    \begin{aligned}
        &d_{s_10}+d_{s_20}+\ldots+d_{s_m0},\\
        \geq&d_{s_11}+d_{s_20}+\ldots+d_{s_m0},\\
        \geq &\ldots,\\
        \geq&d_{s_1l_1}+d_{s_20}+\ldots+d_{s_m0}\\
        \geq &\ldots,\\
        \geq&d_{s_1l_1}+d_{s_2l_2}+\ldots+d_{s_m0}\\
        \geq &\ldots,\\
        \geq&d_{s_1l_1}+d_{s_2l_2}+\ldots+d_{s_ml_m}\\
    \end{aligned}
\end{equation*}
\item Let $S_j$ be the first sum in chain $S$ which was not popped from queue $Q$ before the $j^{th}$ iteration. During the $j^{th}$ iteration, when $S_{j-1}$ is popped, $S_{j}$ is added to $Q$. 
\item Thus, during the $k^{th}$ iteration, the popped sum $S_i$ satisfies $S_i\leq S_j\leq S_l = d_{s_1 l_1}+d_{s_2l_1}+\ldots +d_{s_ml_m}$.
\item By induction, since the $(i-1)^{th}$ largest sum has been popped before the $i^{th}$ iteration,  so the sum popped in the $i^{th}$ iteration is exactly the $i^{th}$ largest sum0 $S_l = d_{s_1 l_1}+d_{s_2l_1}+\ldots +d_{s_ml_m}$.
\end{itemize}
Then we finish the proof.
\end{proof}
A simple example is shown when $m=2$ in Fig.\ref{appe:exact}.

\begin{figure*}[!t]
\centering
\includegraphics[width=2.0\columnwidth]{chap4_exact.png}
\caption{An example of exact evaluation when $m=2$} \label{appe:exact}
\end{figure*}

\section{A Theoretical Quantitative Analysis between Estimation Error and Sample Size}
To theoretically quantify the relationship between estimation error and sample size of data, we construct a new random variable $S$ that is the sum of multiple random variables such that the variable itself is a random variable. 
Then we define a Bernoulli variable $Z\in \{0,1\}$, i.e.,
\begin{equation}
    Z=\left\{
\begin{aligned}
1, &\quad\quad S\leq W \\
0, &\quad\quad S> W
\end{aligned}
\right.
\end{equation}
for a combination of items with $L$ sample size, where $p$ is the true confidence level, i.e., the expectation $\mathbb{E}Z$ of the variable $Z$. 
Hoeffding's inequality shows that for independent identically distributed random variables $Z_1,Z_2,\ldots,Z_L$ taking values $[a,b]$, for any $\epsilon>0$, the following inequality holds,
\begin{equation}
    Pr(|\mathbb{E}[Z]-\frac{1}{L}\sum_{l=1}^L Z_l|\geq \epsilon)\leq 2exp(-\frac{2L\epsilon^2}{(b-a)^2})
\end{equation}
For our problem, we are concerned with the true value is lower than the estimate and greater than the error, i.e., where $\epsilon$ is the estimation error corresponding to the confidence level. 
For the Bernoulli distribution, we have $\mathbb{E}Z=p$ and $a=0,b=1$. 
Thus, the inequality becomes
\begin{equation}
    Pr(p-f(Z_1,Z_2,\ldots,Z_L)\leq -\epsilon)\leq exp(-2L\epsilon^2)
\end{equation}
where $f(Z_1,Z_2,\ldots,Z_L)$ denotes the confidence evaluation algorithm on some combination of items. 
Assuming that at least 5 of the 10 solutions in the output solution set are true feasible solutions, and $f(Z_1,Z_2,\ldots,Z_L)$ is accurate, then we have $1-exp(-2L\epsilon^2)\geq 5/10$. 
Finally, it leads to 
\begin{equation}
    L\geq \frac{ln2}{2\epsilon^2}
\end{equation}

As this bound is derived from Hoeffding's inequality, Chernoff bound is proposed in \cite{tempo1996probabilistic} and a more tight one is proposed in \cite{alamo2010sample}. 
However, all of these inequalities indicate that a large amount of samples is necessary. 
The curve of the variation of $L$ with $\epsilon$ is shown in Figure \ref{fig:samplesize_curve}.

\begin{figure}[!t]
\centering
\includegraphics[width=1.0\columnwidth]{chap4_samplesize_curve.png}
\caption{Sample size requirement with different error in theory} \label{fig:samplesize_curve}
\end{figure}
The minimum sample sizes required for different estimation error are listed in Table \ref{tab:samplesize}. 

Overall, Hoeffding provides the tightest bounds, especially suitable for larger errors. 
However, as error requirements become smaller, the bounds provided by Alamo are more advantageous. 
Practically, it is verified that the amount of data required to meet the corresponding error requirement for the true confidence will be much smaller than the theoretical value obtained.

\begin{table}[!t]
\centering
\caption{Relationship between EE and sample size}
\begin{tabular}{ccccc}
\toprule
\bfseries Estimation error($\%$) & 5 & 0.5 & 0.05 & 0.005 \\
\midrule
{\bfseries Hoeffding's bound} &139 &13863 &1386295 & 138629437\\
{\bfseries Chernoff's bound \cite{tempo1996probabilistic}} &278 &27726 &2772589 &277258873\\
{\bfseries Alamo's bound \cite{alamo2010sample}} &1573 &15723 &157227 &1572269\\
\bottomrule
\label{tab:samplesize}
\end{tabular}
\end{table}


\section{Deterministic method based on Gaussian Assumption and Mix Integer Linear Programming for DDCCMCKP}
\subsection{Risk-averse knapsack problem method}
Yang and Chakraborty \cite{yang2018algorithm} proposed an algorithm for the chance-constrained knapsack problem (CCKP) assuming known Gaussian distributions of the random variables (mean and variance are known). 
CCKP was then transformed into an iterative Risk-averse Knapsack Problem (RA-KP), which could be solved via dynamic programming. 
The implement on DDCCMCKP of this work is easy. 
We first calculated the sample mean and standard deviation of the item weights, $\mu_{ij}=\mathbb{E}(w_{ij}),\sigma^2_{ij}=Var(w_{ij}),i\in \{1,...,m\}, j \in N_i$, and then plug them into the algorithm of Yang's,

\begin{subequations}
    \begin{align}
    & & \min \quad       &\sum\limits_{i=1}^m\sum\limits_{j\in N_i} c_{ij} x_{ij}    \\
    & & \text{s.t.}\quad &\sum\limits_{i=1}^m\sum\limits_{j\in N_i} \mu_{ij} x_{ij} + \Phi^{-1}(P_0)\sqrt{\sum\limits_{i=1}^m\sum\limits_{j\in N_i} \sigma^2_{ij} x_{ij}}\leq W\\
    & &                  &\sum\limits_{j\in N_i}x_{ij}=1, \forall i \in \mathcal{M}\\
    & &                  &x_{ij}\in \{0,1\},\forall i \in \mathcal{M}, j\in N_i 
    \end{align}
    \label{formu_RAMCKP}
\end{subequations}
where $\Phi$ represents the CDF of a standard normal distribution with zero mean and unit variance.

\subsection{Mix Integer Linear Programming}
Ji and Lejeune \cite{ji2021data} studied the distribution-robust chance-constrained planning (DRCCP) optimization problem with data-driven Wasserstein fuzzy sets and proposed an algorithm and reconstruction framework suitable for all types of distribution-robust chance-constrained optimization problems, which was then applied to multidimensional knapsack problems. 
When the decision variables are binary, the proposed LP (who becomes MILP) formulations are equivalent reformulations of the CC. 
To fit the CC into \cite{ji2021data}'s modelling and solution framework, we first formulate $\text{Prob}(\sum\limits_{i=1}^m\sum\limits_{j\in N_i} w_{ij} x_{ij}\leq W)\geq P_0$ as:
\begin{equation}\label{appe:ap.1}
    \text{Prob}(\sum\limits_{i=1}^m\sum\limits_{j\in N_i} w_{ij} x_{ij} > W)\leq 1-P_0
\end{equation}

Define $b=-W, \epsilon = 1-P_0, \tilde{w}_{ij}=-w_{ij}, \forall i \in \mathcal{M}, j\in N_{i}$. 
Let $\bar{x}\in\{0,1\}^{NM}$ as the vectorization of $x\in\{0,1\}^{N\times M}$, $\xi \in\mathbb{R}^{NM}$ is the vectorization of $\hat{w}\in\mathbb{R}^{N\times M}$. 
Finally, we have the standard form of \ref{appe:ap.1} as \cite{ji2021data}:
\begin{equation}\label{appe:ap.3}
    \text{Prob}(\xi^{T}\bar{x}< b)\leq \epsilon
\end{equation}
Then we can utilize Theorem 10 in \cite{ji2021data} to reformulate the above inequation \ref{appe:ap.1}, and get the MILP formulation of DDCCMCKP: 
\begin{equation}
\begin{aligned}
    &(\bar{x},\lambda,\eta,s,\phi)\in \mathcal{Z}_{LP}^{MCKP}\\
    =&\left\{
        \begin{array}{lr}
        \lambda \theta+\frac{1}{L}\sum _{l=1}^L s_l\leq \epsilon, & \\
        \lambda\geq 0, &\\
        s_l\geq 0,\eta_l\geq 0,& \forall l\in \mathcal{L}\\
        1+\eta_lb^{'}-\sum_{q\in\mathcal{Q}} \xi_{ql}^0\phi_{ql}\leq s_l,&\forall l\in \mathcal{L}\\
        \phi_{ql}\leq \lambda, &\forall q\in\mathcal{Q}, l\in \mathcal{L}\\
        \phi_{ql}\leq \eta_l, &\forall q\in\mathcal{Q}, l\in \mathcal{L}\\
        \phi_{ql}\leq U_{\eta_l}\bar{x}_q, &\forall q\in\mathcal{Q}, l\in \mathcal{L}\\
        \phi_{ql}\geq 0, &\forall q\in\mathcal{Q}, l\in \mathcal{L}\\
        \phi_{ql}\geq \eta_l+U_{\eta_{l}}\bar{x}_q -U_{\eta_{l}}, &\forall q\in\mathcal{Q}, l\in \mathcal{L}\\
        \sum _{i=Nm-N+1}^{Nm} \bar{x}_i = 1,&\forall m\in\mathcal{M}\\
        \bar{x}_q \in \{0,1\},&\forall q\in\mathcal{Q}\\
        \end{array}
    \right.
    \label{MILP}
\end{aligned}
\end{equation}

where $b^{'} = b-\delta$, $\delta $ is a user-defined infinite small positive number. 
$\mathcal{Q}$ is the set of all $N\times M$ items, $\mathcal{M}$ is the set of $M$ classes. 
Note that $\sum _{i=Nm-N+1}^{Nm} \bar{x}_i = 1,\forall m\in\mathcal{M}$ is multiple-choice constraint in MCKP. 
$U_{\eta}$ is the upper bound of $\eta$ and can be determined by:
\begin{equation}
    U_{\eta_l}=\left\{
        \begin{array}{lr}
        \mathcal{K}^{'*}, &if \exists \bar{x}: \bar{x}^{T}\xi_l^{0}-b>0\\
        0, &otherwise\\
        \end{array}
    \right.
\end{equation}
where $\mathcal{K}^{'*}$ is the optimal objective value of the linear programming problem:
\begin{equation}
    \begin{aligned}
    \mathcal{K}^{'*} = &\min \bar{x}^{T}\xi_l^{0}-b\\
    \text{s.t.} & b-\bar{x}^{T}\xi_l^{0}\leq 0-\tau ^{'}\\
    & \bar{x} \in \mathcal{\bar{X}}
    \end{aligned}
    \label{appe:LP for U}
\end{equation}
where $\tau^{'}$ a user-defined small non-negative number.

Overall, we first solve linear programming \ref{appe:LP for U} to get $U_{\eta_l},\forall l \in \mathcal{L}$ then solve MILP $\mathcal{Z}_{LP}^{MCKP}$ \ref{MILP} to find the optimal solution of DDCCMCKP.

\section{Complete experimental results}
In this section, we show the complete experimental results of DDALS on LAB and APP.

\begin{table*}[!t]
    \caption{Performance on benchmarks of LAB. The highest average real confidence level and the lowest cost are highlighted through bolding and being underlined.}
    \centering
    \resizebox{\linewidth}{!}{
    \begin{tabular}{ccccccccccc}
    \toprule
    \multirowcell{2}{\bfseries Benchmark} & \multicolumn{5}{c}{\bfseries small $T_{\max}$} & \multicolumn{5}{c}{\bfseries large $T_{\max}$} \\
    \cmidrule(r){2-6} \cmidrule(r){7-11}
    & {\bfseries C} & {\bfseries ET} & {\bfseries ECL} & {\bfseries RCL} & {\bfseries FSR} & {\bfseries C} & {\bfseries ET} & {\bfseries ECL} & {\bfseries RCL} & {\bfseries FSR}\\
    \midrule
    {\bfseries ss1-Greedy} &22.2$\pm$0.0 &2.0$\pm$0.0 & 0$\pm$0.0 &99.6$\pm$0.0 &\underline{\bfseries 1} &15.1$\pm$0.0 &3.0$\pm$0.0& 99.9$\pm$0.0 &99.9$\pm$0.0&\underline{\bfseries 1} \\
    {\bfseries ss1-GA}  &18.7$\pm$0.1 &354$\pm$4.2 & 99.8$\pm$0.3 &99.6$\pm$0.8 &0.8 &\underline{\bfseries10.8$\pm$0.0} &253.3$\pm$4.1& 99.9$\pm$0.1 &99.0$\pm$1.2  &0.6 \\
    {\bfseries ss1-EDA} &21.5$\pm$1.4 &354.5$\pm$2.9 & - &99.5$\pm$0.5 &0.1 &13.0$\pm$3.4 &234.7$\pm$2.8& 99.8$\pm$0.3 &99.9$\pm$0.1 &\underline{\bfseries 1} \\
    \bfseries ss1-DDMLS(O) &\underline{\bfseries 18.6$\pm$0.0} &351.7$\pm$5.8 & 99.9$\pm$0.0 &99.9$\pm$0.0 &\underline{\bfseries 1} &\underline{\bfseries 10.8$\pm$0.0}& 231$\pm$13.7 &99.9$\pm$0.0 &99.9$\pm$0.0 &\underline{\bfseries 1} \\ 
    \bfseries ss1-DDMLS(V1) &\underline{\bfseries 18.6$\pm$0.0}  &351.7$\pm$5.8 & 99.9$\pm$0.0 &99.9$\pm$0.0 &\underline{\bfseries 1} &\underline{\bfseries 10.8$\pm$0.0}& 231$\pm$13.7 &99.9$\pm$0.0 &99.9$\pm$0.0 &\underline{\bfseries 1} \\ 
    \bfseries ss1-DDMLS(V2) &\underline{\bfseries 18.6$\pm$0.0}  &351.7$\pm$5.8 & 99.9$\pm$0.0 &99.9$\pm$0.0 &\underline{\bfseries 1} &\underline{\bfseries 10.8$\pm$0.0}& 231$\pm$13.7 &99.9$\pm$0.0 &99.9$\pm$0.0 &\underline{\bfseries 1} \\ 
    {\bfseries ss1-DDMLS(V3)} &20.7$\pm$2.0 &360 & 99.6$\pm$0.3 &99.2$\pm$0.9 &0.3 &13.0$\pm$1.7& 218 &99.8$\pm$0.3 &99.1$\pm$1.0 &0.7 \\ 
    
    \midrule
    {\bfseries ss2-Greedy} &\underline{\bfseries 19.4$\pm$0.0} &2.0$\pm$0.0 & 99.9$\pm$0.0 &99.9$\pm$0.0 &\underline{\bfseries 1} &16.5$\pm$0.0 &4.0$\pm$0.0& 99.9$\pm$0.0 &99.9$\pm$0.0 &\underline{\bfseries 1} \\
    {\bfseries ss2-GA} &\underline{\bfseries 19.4$\pm$0.0} &384.5$\pm$2.6 & 99.9$\pm$0.0 &99.9$\pm$0.0 &\underline{\bfseries 1} &11.1$\pm$1.4 &303.6$\pm$3.4& 99.7$\pm$0.2 &99.7$\pm$0.4 &\underline{\bfseries 1} \\
    {\bfseries ss2-EDA} &19.9$\pm$0.4 &385.9$\pm$2.7 & 99.9$\pm$0.0 &99.9$\pm$0.1 &\underline{\bfseries 1} &14.8$\pm$2.2 &304.5$\pm$2.6& 99.7$\pm$0.2 &99.7$\pm$0.4 &\underline{\bfseries 1} \\
    \bfseries ss2-DDMLS(O) &\underline{\bfseries 19.4$\pm$0.0} &378.8$\pm$8.2 & 99.9$\pm$0.0 &99.9$\pm$0.0 &\underline{\bfseries 1} &\underline{\bfseries10.4$\pm$0.0} &298.9$\pm$15.5& 99.8$\pm$0.0 &99.9$\pm$0.0 &\underline{\bfseries 1} \\
   \bfseries ss2-DDMLS(V1) &\underline{\bfseries 19.4$\pm$0.0} &378.8$\pm$8.2 & 99.9$\pm$0.0 &99.9$\pm$0.0 &\underline{\bfseries 1} &\underline{\bfseries10.4$\pm$0.0} &298.9$\pm$15.5& 99.8$\pm$0.0 &99.9$\pm$0.0 &\underline{\bfseries 1} \\
    {\bfseries ss2-DDMLS(V2)} &\underline{\bfseries19.4$\pm$0.0} &378.8$\pm$8.2 & 99.9$\pm$0.0 &99.9$\pm$0.0 &\underline{\bfseries 1} &12.3$\pm$2.0 &298.9$\pm$15.5& 99.9$\pm$0.1 &99.9$\pm$0.0 &\underline{\bfseries 1} \\
    {\bfseries ss2-DDMLS(V3)} &21.8$\pm$2.4 &377 & 99.9$\pm$0.0 &99.9$\pm$0 &\underline{\bfseries 1} &13.8$\pm$1.6& 319 &99.8$\pm$0.3 &99.9$\pm$0.1 &\underline{\bfseries 1} \\

    \midrule
    {\bfseries ss3-Greedy} &24.4$\pm$0.0&2.0$\pm$0.0 & 100.0$\pm$0.0&99.9$\pm$0.0 &\underline{\bfseries 1} &21.8$\pm$0.0&3.0$\pm$0.0& 100$\pm$0&99.9$\pm$0.0 &\underline{\bfseries 1} \\
    {\bfseries ss3-GA} &22.6$\pm$1.5&303.4$\pm$2.8 & 99.9$\pm$0.0 & 99.7$\pm$0.2 &\underline{\bfseries 1} &12.9$\pm$0.7 & 262.4$\pm$1.6&99.9$\pm$0.0&99.9$\pm$0.1 &\underline{\bfseries 1} \\
    {\bfseries ss3-EDA} &24.1$\pm$2.3&304.1$\pm$2.3 & 99.9$\pm$0.1&99.6$\pm$0.6&0.9&14.0$\pm$2.0&264.5$\pm$2.9& 99.9$\pm$0.0&99.9$\pm$0.1 &\underline{\bfseries 1} \\
    {\bfseries ss3-DDALS(O)} &\underline{\bfseries 21.4$\pm$0.0} &291.9$\pm$6.2 & 99.9$\pm$0.0&99.5$\pm$0.0 &\underline{\bfseries 1} &\underline{\bfseries 12.6$\pm$0.0} &254.4$\pm$4.7& 100.0$\pm$0.0 &99.9$\pm$0.0  &\underline{\bfseries 1} \\
    {\bfseries ss3-DDALS(V1)} &\underline{\bfseries 21.4$\pm$0.0} &291.9$\pm$6.2 & 99.9$\pm$0.0&99.5$\pm$0.0 &\underline{\bfseries 1} &\underline{\bfseries 12.6$\pm$0.0}&254.4$\pm$4.7& 100.0$\pm$0.0&99.9$\pm$0.0  &\underline{\bfseries 1} \\ 
    {\bfseries ss3-DDALS(V2)} &\underline{\bfseries 21.4$\pm$0.0} &291.9$\pm$6.2 & 99.9$\pm$0.0&99.5$\pm$0.0 &\underline{\bfseries 1} &\underline{\bfseries 12.6$\pm$0.0}&254.4$\pm$4.7& 100.0$\pm$0.0&99.9$\pm$0.0  &\underline{\bfseries 1} \\
    {\bfseries ss3-DDALS(V3)} &23.9$\pm$1.7&279 & 99.8$\pm$0.2&99.1$\pm$0.8&0.8&15.7$\pm$1.9&260& 99.9$\pm$0.3&99.6$\pm$0.9  &0.9 \\
    
    \midrule
    {\bfseries ss4-Greedy} &30.1$\pm$0.0&2.0$\pm$0.0 & 99.9$\pm$0.0&99.9$\pm$0.0 &\underline{\bfseries 1} &27.5$\pm$0.0&3.0$\pm$0.0& 99.7$\pm$0.0&99.4$\pm$0.0  &\underline{\bfseries 1} \\
    {\bfseries ss4-GA} &23.8$\pm$5.5&768.9$\pm$8.5 & 99.9$\pm$0.3&99.7$\pm$0.4 &\underline{\bfseries 1} &17.2$\pm$3.0&562.6$\pm$5.7& 99.7$\pm$0.3&99.6$\pm$0.3  &\underline{\bfseries 1} \\
    {\bfseries ss4-EDA} &\underline{\bfseries 17.4$\pm$0.0} &763.4$\pm$2.5 & 99.9$\pm$0.0&99.9$\pm$0.0 &\underline{\bfseries 1} &13.0$\pm$1.0&553.4$\pm$2.7& 99.8$\pm$0.2&98.9$\pm$0.4  &\underline{\bfseries 1} \\
    {\bfseries ss4-DDALS(O)} &\underline{\bfseries 17.4$\pm$0.0} &759.9$\pm$25.8 & 99.9$\pm$0.0&99.9$\pm$0.0 &\underline{\bfseries 1} &\underline{\bfseries 11.2$\pm$0.0} &525.3$\pm$39.1& 99.2$\pm$0.0&99.0$\pm$0.0  &\underline{\bfseries 1} \\
    {\bfseries ss4-DDALS(V1)} &\underline{\bfseries 17.4$\pm$0.0} &759.9$\pm$25.8 & 99.9$\pm$0.0&99.9$\pm$0.0 &\underline{\bfseries 1} &12.1$\pm$0.0&525.3$\pm$39.1& 99.5$\pm$0.0&98.9$\pm$0.0  &\underline{\bfseries 1} \\
    {\bfseries ss4-DDALS(V2)} &\underline{\bfseries 17.4$\pm$0.0} &759.9$\pm$25.8 & 99.9$\pm$0.0&99.9$\pm$0.0 &\underline{\bfseries 1} &14.0$\pm$1.2&525.3$\pm$39.1& 99.9$\pm$0.2&99.0$\pm$0.3  &\underline{\bfseries 1} \\
    {\bfseries ss4-DDALS(V3)} &24.6$\pm$4.0&735 & 99.8$\pm$0.2&99.5$\pm$0.5 &\underline{\bfseries 1} &13.8$\pm$1.3&561& 99.8$\pm$0.3&99.5$\pm$0.5  &\underline{\bfseries 1} \\
    
    \midrule
    {\bfseries ls1-Greedy} &58.8$\pm$0.0 &2.0$\pm$0.0 & 99.8$\pm$0.0 &99.8$\pm$0.0 &\underline{\bfseries 1} &56.0$\pm$0.0 &3.0$\pm$0.0& 99.6$\pm$0.0 &99.5$\pm$0.0 &\underline{\bfseries 1} \\
    {\bfseries ls1-GA} &50.6$\pm$0.9 &3007.2$\pm$4.1 & 99.1$\pm$0.1 &99.2$\pm$0.1 &\underline{\bfseries 1} &41.0$\pm$4.7 &1412.6$\pm$8.3& 99.3$\pm$0.3 &99.2$\pm$0.4 &\underline{\bfseries 1} \\
    {\bfseries ls1-EDA} &54.7$\pm$2.1 &3003.1$\pm$2.2 & 99.2$\pm$0.1 &99.3$\pm$0.1 &\underline{\bfseries 1} &37.7$\pm$0 &1404.6$\pm$2.3& 99.3$\pm$0.0 &99.3$\pm$0.0 &\underline{\bfseries 1} \\
    \bfseries ls1-DDMLS(O) &\underline{\bfseries 49.8$\pm$0.1} &2928.8$\pm$71.5 & 99.3$\pm$0.0 &99.2$\pm$0.0 &\underline{\bfseries 1} &\underline{\bfseries 32.0$\pm$0} &1371.3$\pm$12.4& 99.2$\pm$0.0 &99.1$\pm$0.0 &\underline{\bfseries 1} \\
    {\bfseries ls1-DDMLS(V1)} &52.2$\pm$0.3 &2928.8$\pm$71.5 & 99.6$\pm$0.1 &99.6$\pm$0.1 &\underline{\bfseries 1} &34.6$\pm$0.9 &1371.3$\pm$12.4& 99.6$\pm$0.0 &99.6$\pm$0.1 &\underline{\bfseries 1} \\
    {\bfseries ls1-DDMLS(V2)} &49.9$\pm$0.3 &2928.8$\pm$71.5 & 99.3$\pm$0.0 &99.3$\pm$0.1 &\underline{\bfseries 1} &33.1$\pm$1.1 &1371.3$\pm$12.4& 99.4$\pm$0.2 &99.3$\pm$0.0 &\underline{\bfseries 1} \\
    {\bfseries ls1-DDMLS(V3)} &51.2$\pm$0.9 &2919 & 99.3$\pm$0.2 &99.3$\pm$0.3 &\underline{\bfseries 1} &35.1$\pm$1.5 &1383& 99.3$\pm$0.2 &99.2$\pm$0.3 &\underline{\bfseries 1} \\
    
    \midrule
    {\bfseries ls2-Greedy} &66.6$\pm$0.0 &2.0$\pm$0.0 & 99.7$\pm$0.0 &99.7$\pm$0.0 &\underline{\bfseries 1} &66.6$\pm$0.0 &2.0$\pm$0.0& 99.9$\pm$0.0 &100.0$\pm$0.0 &\underline{\bfseries 1} \\
    {\bfseries ls2-GA} &53.1$\pm$2.9 &9623.1$\pm$21.8 & 99.2$\pm$0.1 &99.2$\pm$0.2 &\underline{\bfseries 1} &43.3$\pm$3.1 &3317.4$\pm$ 17.2& 99.6$\pm$0.3 &99.5$\pm$0.3 &\underline{\bfseries 1} \\
    {\bfseries ls2-EDA} &53.5$\pm$3.1 &9604.4$\pm$2.4 & 99.1$\pm$0.0 &99.2$\pm$0.0 &\underline{\bfseries 1} &38.7$\pm$2.7 &3304.6$\pm$ 2.5& 99.5$\pm$0.3 &99.5$\pm$0.2 &\underline{\bfseries 1} \\
    \bfseries ls2-DDMLS(O) &\underline{\bfseries 47.9$\pm$0.0} &9548.9$\pm$46.4 & 99.0$\pm$0.0 &99.0$\pm$0.0 &\underline{\bfseries 1} &\underline{\bfseries 27.2$\pm$0.0} &3305.6$\pm$342.6& 99.3$\pm$0.0 &99.3$\pm$0.0 &\underline{\bfseries 1} \\
    {\bfseries ls2-DDMLS(V1)} &51.3$\pm$0.0 &9548.9$\pm$46.4 & 99.4$\pm$0.0 &99.5$\pm$0.0 &\underline{\bfseries 1} &27.9$\pm$0.4 &3305.6$\pm$342.6& 99.6$\pm$0.0 &99.6$\pm$0.0 &\underline{\bfseries 1} \\
    {\bfseries ls2-DDMLS(V2)} &49.1$\pm$0.0 &9548.9$\pm$46.4 & 99.2$\pm$0.0 &99.2$\pm$0.0 &\underline{\bfseries 1} &28.0$\pm$0.5 &3305.6$\pm$342.6& 99.6$\pm$0.1 &99.6$\pm$0.1 &\underline{\bfseries 1} \\
    {\bfseries ls2-DDMLS(V3)} &51.2$\pm$2.4 &9518 & 99.2$\pm$0.1 &99.2$\pm$0.1 &\underline{\bfseries 1} &29.2$\pm$1 &3197& 99.4$\pm$0.2 &99.4$\pm$0.3 &\underline{\bfseries 1} \\
    
    \midrule
    {\bfseries ls3-Greedy} &101.8$\pm$0.0 &3.0$\pm$0.0 & 99.9$\pm$0.0 &99.9$\pm$0.0 &\underline{\bfseries 1} &101.8$\pm$0.0 &4.0$\pm$0.0& 99.8$\pm$0.0 &99.9$\pm$0.0 &\underline{\bfseries 1} \\
    {\bfseries ls3-GA} &85.7$\pm$5.6 &7213.2$\pm$11.3 & 99.1$\pm$0.1 &98.8$\pm$0.1 &\underline{\bfseries 1} &71.3$\pm$4.0 &5209.5$\pm$7.2& 99.4$\pm$0.3 &99.3$\pm$0.3 &\underline{\bfseries 1} \\
    {\bfseries ls3-EDA} &82.5$\pm$0.4 &7202.6$\pm$2.8 & 99.1$\pm$0.2 &98.8$\pm$0.2 &\underline{\bfseries 1} &73.4$\pm$3.3 &5203.7$\pm$2.8& 99.7$\pm$0.3 &99.6$\pm$0.4 &\underline{\bfseries 1} \\
    \bfseries ls3-DDMLS(O) &\underline{\bfseries 74.9$\pm$0.2} &7233.5$\pm$58.3 & 99.1$\pm$0.0 &98.7$\pm$0.0 &\underline{\bfseries 1} &\underline{\bfseries 63.4$\pm$0.9} &5217.5$\pm$274.1& 99.3$\pm$0.2 &99.1$\pm$0.2 &\underline{\bfseries 1} \\
    {\bfseries ls3-DDMLS(V1)} &79.6$\pm$1.3 &7233.5$\pm$58.3 & 99.6$\pm$0.1 &99.4$\pm$0.1 &\underline{\bfseries 1} &65.2$\pm$2.1 &5217.5$\pm$274.1& 99.6$\pm$0.1 &99.5$\pm$0.2 &\underline{\bfseries 1} \\
    {\bfseries ls3-DDMLS(V2)} &75.9$\pm$0.7 &7233.5$\pm$58.3 & 99.4$\pm$0.1 &99.1$\pm$0.1 &\underline{\bfseries 1} &65.0$\pm$1.9 &5217.5$\pm$274.1& 99.7$\pm$0.2 &99.6$\pm$0.2 &\underline{\bfseries 1} \\
    {\bfseries ls3-DDMLS(V3)} &77.3$\pm$1.5 &7175 & 99.3$\pm$0.1 &99.0$\pm$0.2 &\underline{\bfseries 1} &65.2$\pm$0.8 &5109& 99.4$\pm$0.3 &99.3$\pm$0.4 &\underline{\bfseries 1} \\

    \midrule
    {\bfseries ls4-Greedy} &169.7$\pm$0.0&5$\pm$0.0 & 99.8$\pm$0.0&99.8$\pm$0.0 &\underline{\bfseries 1} &162.5$\pm$0.0&6.0$\pm$0.0& 99.9$\pm$0.0&99.9$\pm$0.0  &\underline{\bfseries 1} \\
    {\bfseries ls4-GA} &123.1$\pm$8.8&11307.2$\pm$8.0 & 99.1$\pm$0.1&99.1$\pm$0.2 &\underline{\bfseries 1} &104.1$\pm$5.1&8409.1$\pm$5.6& 99.1$\pm$0.2&99.2$\pm$0.2  &\underline{\bfseries 1} \\
    {\bfseries ls4-EDA} &126.2$\pm$2.3&11304.8$\pm$2.1 & 99.2$\pm$0.2&99.4$\pm$0.2 &\underline{\bfseries 1} &114.8$\pm$4.5& 8403.4$\pm$2.8&99.6$\pm$0.2&99.6$\pm$0.2  &\underline{\bfseries 1} \\
    {\bfseries ls4-DDALS(O)} &\underline{\bfseries 106.6$\pm$0.8} &11289.9$\pm$92.0 & 99.1$\pm$0.1&99.2$\pm$0.1 &\underline{\bfseries 1} &\underline{\bfseries 92.8$\pm$0.7} &8376.6$\pm$378.8& 99.1$\pm$0.1&99.2$\pm$0.1  &\underline{\bfseries 1} \\
    {\bfseries ls4-DDALS(V1)} &109.5$\pm$0.4&11289.9$\pm$92.0 & 99.7$\pm$0.1&99.7$\pm$0.1 &\underline{\bfseries 1} &97.9$\pm$2.2&8376.6$\pm$378.8& 99.6$\pm$0.1&99.7$\pm$0.1  &\underline{\bfseries 1} \\
    {\bfseries ls4-DDALS(V2)} &108.9$\pm$1.3&11289.9$\pm$92.0 & 99.6$\pm$0.2&99.7$\pm$0.2 &\underline{\bfseries 1} &93.8$\pm$1.4&8376.6$\pm$378.8& 99.3$\pm$0.2&99.4$\pm$0.2  &\underline{\bfseries 1} \\
    {\bfseries ls4-DDALS(V3)} &108.7$\pm$1.1&11171 & 99.2$\pm$0.2&99.3$\pm$0.2 &\underline{\bfseries 1} &96.5$\pm$1.7&8177& 99.1$\pm$0.1& 99.2$\pm$0.1  &\underline{\bfseries 1} \\
    
    \midrule
    {\bfseries ls5-Greedy} &213.8$\pm$0.0 &3.0$\pm$0.0 & 99.9$\pm$0.0 &99.8$\pm$0.0 &\underline{\bfseries 1} &205.5$\pm$0.0 &4.0$\pm$0.0& 99.9$\pm$0.0 &99.9$\pm$0.0 &\underline{\bfseries 1} \\
    {\bfseries ls5-GA} &192.6$\pm$6.7 &29013.4$\pm$7.6 & 99.0$\pm$0.0 &98.8$\pm$0.0 &\underline{\bfseries 1} &170.2$\pm$7.7 &18313$\pm$10.3& 99.1$\pm$0.2 &99.0$\pm$0.3 &0.9 \\
    {\bfseries ls5-EDA} &195.2$\pm$2.2 &29004.2$\pm$3.1 & 99.1$\pm$0.1 &99.0$\pm$0.2 &\underline{\bfseries 1} &165.9$\pm$1.9 &18305.5$\pm$2.8& 99.3$\pm$0.2 &99.3$\pm$0.2 &\underline{\bfseries 1} \\
    \bfseries ls5-DDMLS(O) &\underline{\bfseries 178.1$\pm$0.8} &28840.1$\pm$289.5 & 99.0$\pm$0.0 &98.9$\pm$0.0 &\underline{\bfseries 1} &\underline{\bfseries 144.5$\pm$1.0} &18276.1$\pm$395.4& 99.0$\pm$0.0 &99.0$\pm$0.1 &\underline{\bfseries 1} \\
    {\bfseries ls5-DDMLS(V1)} &184.5$\pm$1.4 &28840.1$\pm$289.5 & 99.3$\pm$0.0 &99.2$\pm$0.0 &\underline{\bfseries 1} &155.8$\pm$3.3 &18276.1$\pm$395.4& 99.6$\pm$0.1 &99.5$\pm$0.1 &\underline{\bfseries 1} \\
    {\bfseries ls5-DDMLS(V2)} &179.6$\pm$1.9 &28840.1$\pm$289.5 & 99.2$\pm$0.0 &99.0$\pm$0.0 &\underline{\bfseries 1} &147.3$\pm$3.4 &18276.1$\pm$395.4& 99.2$\pm$0.2 &99.1$\pm$0.2 &\underline{\bfseries 1} \\
    {\bfseries ls5-DDMLS(V3)} &182.9$\pm$1.6 &28624 & 99.1$\pm$0.0 &99.0$\pm$0.0 &\underline{\bfseries 1} &147.4$\pm$4.0 &17573& 99.1$\pm$0.2 &99.1$\pm$0.2 &\underline{\bfseries 1} \\
    
    \midrule
    {\bfseries ls6-Greedy} &256.3$\pm$0.0 &4.0$\pm$0.0 & 99.4$\pm$0.0 &99.4$\pm$0.0 &\underline{\bfseries 1} &250.0$\pm$0.0 &7.0$\pm$0.0& 99.8$\pm$0.0 &99.7$\pm$0.0 &\underline{\bfseries 1} \\
    {\bfseries ls6-GA} &233.5$\pm$14.0 &34609.3$\pm$6.0 & 99.0$\pm$0.0 &98.9$\pm$0.0 &\underline{\bfseries 1} &187.2$\pm$9.3 &20211.9$\pm$8.3& 99.0$\pm$0.0 &99.0$\pm$0.2 &0.9 \\
    {\bfseries ls6-EDA} &214.1$\pm$2.8 &34604.9$\pm$2.8 & 99.3$\pm$0.2 &99.2$\pm$0.2 &\underline{\bfseries 1} &181.7$\pm$1.8 &20205.1$\pm$2.7& 99.4$\pm$0.2 &99.3$\pm$0.2 &\underline{\bfseries 1} \\
    \bfseries ls6-DDMLS(O) &\underline{\bfseries 198.5$\pm$ 0.5} & 34570.7$\pm$619.3 &99.0$\pm$0.0 &98.9$\pm$0.0 &\underline{\bfseries 1} &\underline{\bfseries 162.2$\pm$1.1} &20178.8$\pm$1031.6 &99.1$\pm$0.1& 99.0$\pm$0.1 &\underline{\bfseries }1 \\
    {\bfseries ls6-DDMLS(V1)} &206.8$\pm$ 2.3 & 34570.7$\pm$619.3 &99.5$\pm$0.0 &99.4$\pm$0.1 &\underline{\bfseries 1}  &169.6$\pm$1.4 &20178.8$\pm$1031.6 &99.5$\pm$0.0& 99.5$\pm$0.1 &\underline{\bfseries 1} \\
    {\bfseries ls6-DDMLS(V2)} &201.5$\pm$ 2.4 & 34570.7$\pm$619.3 &99.3$\pm$0.2 &99.2$\pm$0.2 &\underline{\bfseries 1}  &163.1$\pm$1.5 &20178.8$\pm$1031.6 &99.3$\pm$0.1& 99.3$\pm$0.1 &\underline{\bfseries 1} \\
    {\bfseries ls6-DDMLS(V3)} &202.6$\pm$ 2.0 & 33681 &99.1$\pm$0.1 &99.1$\pm$0.1 &\underline{\bfseries 1}  &165.0$\pm$2.6 &19387 &99.2$\pm$0.2& 99.1$\pm$0.1 &\underline{\bfseries 1} \\

    \bottomrule
    \end{tabular}
    }
\label{fullresult:LAB}
\end{table*}

\begin{table*}[!t]
    \caption{Performance on benchmarks of APP. Due to the introduction of retransmission mechanisms, there is a significant difference in the results between LAB and APP.}
    \centering
    \resizebox{\linewidth}{!}{
        \begin{tabular}{ccccccccccc}
        \toprule
        \multirowcell{2}{\bfseries Benchmark} & \multicolumn{5}{c}{\bfseries small $T_{\max}$} & \multicolumn{5}{c}{\bfseries large $T_{\max}$} \\
        \cmidrule(r){2-6} \cmidrule(r){7-11}
        & {\bfseries C} & {\bfseries ET} & {\bfseries ECL} & {\bfseries RCL} & {\bfseries FSR} & {\bfseries C} & {\bfseries ET} & {\bfseries ECL} & {\bfseries RCL} & {\bfseries FSR}\\
        
        \midrule
        {\bfseries ss1-Greedy} &9.2$\pm$0.0 &4.0$\pm$0.0 & 99.3$\pm$0.0 &96.7$\pm$0 &0 &9.0$\pm$0.0 &60.0$\pm$0.0& 99.6$\pm$0.0 &98.3$\pm$0.0 &0 \\
        {\bfseries ss1-GA} &10.5$\pm$2.3 &337.4$\pm$5.2 & 99.4$\pm$0.3 &96.9$\pm$0.7 &0 &8.6$\pm$2.1 &303.7$\pm$2.7& 99.5$\pm$0.1 &97.2$\pm$0.0 &0 \\
        {\bfseries ss1-EDA} &\underline{\bfseries9.0$\pm$0.0} &332.5$\pm$2.3 & 99.3$\pm$0.0 &96.5$\pm$0.0 &0 &8.0$\pm$0.9 &304.2$\pm$2.4& 99.6$\pm$0.0 &97.7$\pm$0.6 &0 \\
        {\bfseries ss1-DDMLS(O)} &\underline{\bfseries9.0$\pm$0.0} &311.4$\pm$10.2 & 99.2$\pm$0.0 &96.5$\pm$0.0 &0 &\underline{\bfseries7.3$\pm$0.0} &285.5$\pm$14.8& 99.6$\pm$0.0 &97.2$\pm$0.0 &0 \\
        {\bfseries ss1-DDMLS(V1)} &14.0$\pm$0.0 &311.4$\pm$10.2 & 99.9$\pm$0.0 &97.2$\pm$0.0 &0 &\underline{\bfseries7.3$\pm$0.0} &285.5$\pm$14.8& 99.6$\pm$0.0 &97.2$\pm$0.0 &0 \\
        {\bfseries ss1-DDMLS(V2)} &13.5$\pm$1.5 &311.4$\pm$10.2 & 99.8$\pm$0.2 &97.1$\pm$0.2 &0 &\underline{\bfseries7.3$\pm$0.0} &285.5$\pm$14.8& 99.6$\pm$0.0 &97.2$\pm$0.0 &0\\
        {\bfseries ss1-DDMLS(V3)} &14.8$\pm$3.4 &321 & 99.7$\pm$0.3 &98.0$\pm$1.0 &\underline{\bfseries0.5} &11.1$\pm$2.4 &300& 99.5$\pm$0.0 &97.7$\pm$0.6 &\underline{\bfseries0.1} \\
        
        \midrule
        {\bfseries ss2-Greedy} & 6.4$\pm$0.0 & 3$\pm$0.0 & 99.7$\pm$0.0 &98.6$\pm$0.0 &\underline{\bfseries 1} &6.4$\pm$0.0&3.0$\pm$0.0& 99.8$\pm$0.0&99.0$\pm$0.0  &\underline{\bfseries 1} \\
        {\bfseries ss2-GA} &5.3$\pm$0.0&301.8$\pm$2.2 & 99.2$\pm$0.0&97.3$\pm$0.4&0&5.3$\pm$0.0&272.6$\pm$2.8& 99.5$\pm$0.1&98.1$\pm$0.2  &0 \\
        {\bfseries ss2-EDA} &5.4$\pm$0.2&303.1$\pm$3.0 & 99.3$\pm$0.2&97.4$\pm$0.7&0.1&5.4$\pm$0.2&274.9$\pm$3.0& 99.6$\pm$0.2&98.3$\pm$0.4  &0.3 \\
        {\bfseries ss2-DDALS(O)} &\underline{\bfseries 5.2$\pm$0.0} &307$\pm$7.6 & 99.2$\pm$0.0&97.3$\pm$0.0&0 &\underline{\bfseries 5.2$\pm$0.0} &291$\pm$10.6& 99.6$\pm$0.0&98.1$\pm$0.0  &0 \\
        {\bfseries ss2-DDALS(V1)} &5.5$\pm$0.1&307$\pm$7.6 & 99.8$\pm$0.0&97.5$\pm$0.3&0 &\underline{\bfseries 5.2$\pm$0.0} &291$\pm$10.1& 99.6$\pm$0.0&98.1$\pm$0.0  &0 \\
        {\bfseries ss2-DDALS(V2)} &5.5$\pm$0.1&307$\pm$7.6 & 99.8$\pm$0.0&97.5$\pm$0.3&0&5.4$\pm$0.1&291$\pm$10.1& 99.9$\pm$0.1&98.2$\pm$0.1  &0 \\
        {\bfseries ss2-DDALS(V3)} &5.4$\pm$0.1&295 & 99.4$\pm$0.3&97.7$\pm$0.4&0&5.3$\pm$0.1&272.6$\pm$2.8& 99.5$\pm$0.1&98.1$\pm$0.2  &0 \\
        
        \midrule
        {\bfseries ss3-Greedy} &18.3$\pm$0.0 &6.0$\pm$0.0 & 99.9$\pm$0.0 &96.3$\pm$0 &0 &13.8$\pm$0.0 &8.0$\pm$0.0 & 99.6$\pm$0.0 &97.6$\pm$0.0 &0 \\
        {\bfseries ss3-GA} &17.5$\pm$0.3 &634.6$\pm$4.0 & 99.3$\pm$0.3 &95.3$\pm$0.5 &0 &12.1$\pm$1.0 &503.6$\pm$2.5& 99.2$\pm$0.2 &94.0$\pm$1.6 &0 \\
        {\bfseries ss3-EDA} &\underline{\bfseries17.1$\pm$0.3} &634.5$\pm$3.5 & 99.1$\pm$0.2 &93.9$\pm$1.1 &0 &11.2$\pm$0.2 &503.7$\pm$3.1& 99.5$\pm$0.2 &94.7$\pm$0.8 &0 \\
        {\bfseries ss3-DDMLS(O)} &\underline{\bfseries17.1$\pm$0.2} &675.9$\pm$36.2 & 99.1$\pm$0.1 &94.2$\pm$0.7 &0 &\underline{\bfseries11.1$\pm$0.1} &499.7$\pm$21.7& 99.5$\pm$0.2 &94.6$\pm$0.7 &0 \\
        {\bfseries ss3-DDMLS(V1)} &18.1$\pm$0.2 &675.9$\pm$36.2 & 99.6$\pm$0.1 &96.7$\pm$0.6 &0 &11.5$\pm$1.0 &499.7$\pm$ 1.7& 99.6$\pm$0.0 &95.4$\pm$1.1 &0 \\
        {\bfseries ss3-DDMLS(V2)} &18.1$\pm$0.3 &675.9$\pm$36.2 & 99.8$\pm$0.2 &96.7$\pm$0.7 &0 &11.3$\pm$0.8 &499.7$\pm$21.7& 99.6$\pm$0.2 &94.9$\pm$1.0 &0 \\
        {\bfseries ss3-DDMLS(V3)} &19.0$\pm$0.8 &632 & 99.5$\pm$0.3 &96.3$\pm$0.5 &0 &12.4$\pm$0.8 &499& 99.5$\pm$0.2 &95.5$\pm$0.1 &0 \\
        
        \midrule
        {\bfseries ss4-Greedy} &34.9$\pm$0.0 &3.0$\pm$0.0 & 100.0$\pm$0.0&100.0$\pm$0.0 &\underline{\bfseries 1} &27.2$\pm$0.0&4.0$\pm$0.0& 100.0$\pm$0.0&99.0$\pm$0.0  &\underline{\bfseries 1} \\
        {\bfseries ss4-GA} &27.3$\pm$1.1&2010$\pm$8.4 & 99.8$\pm$0.3&90.4$\pm$0.5&0&21.5$\pm$2.9&1614.8$\pm$9.8& 99.6$\pm$0.3&93.2$\pm$4.8  &0 \\
        {\bfseries ss4-EDA} &\underline{\bfseries 26.1$\pm$0.0} &2005.7$\pm$2.3 & 100.0$\pm$0.0&90.1$\pm$0.0&0 &\underline{\bfseries 16.4$\pm$0.0} &1602.7$\pm$2.1& 99.9$\pm$0.0&80.3$\pm$0.0  &0 \\
        {\bfseries ss4-DDALS(O)} &\underline{\bfseries 26.1$\pm$0.0} &1988.2$\pm$30.4 & 100.0$\pm$0.0&90.1$\pm$0.0&0&16.7$\pm$0.9&1579.4$\pm$76.8 & 99.8$\pm$0.2&81.8$\pm$4.4&0 \\
        {\bfseries ss4-DDALS(V1)} &\underline{\bfseries 26.1$\pm$0.0} &1988.2$\pm$30.4 & 100.0$\pm$0.0&90.1$\pm$0.0&0&16.7$\pm$1.0&1579.4$\pm$76.8& 99.9$\pm$0.0&81.5$\pm$3.6  &0 \\
        {\bfseries ss4-DDALS(V2)} &\underline{\bfseries 26.1$\pm$0.0} &1988.2$\pm$30.4 & 100.0$\pm$0.0&90.1$\pm$0.0&0&18.0$\pm$1.9&1579.4$\pm$76.8& 99.9$\pm$0.1&86.1$\pm$7.0  &0 \\
        {\bfseries ss4-DDALS(V3)} &29.1$\pm$2.4&2022 & 99.9$\pm$0.2&91.3$\pm$2.9&0&19.9$\pm$1.7&1589& 99.7$\pm$0.3&92.0$\pm$5.3  &0 \\
        
        \midrule
        {\bfseries ls1-Greedy} &84.1$\pm$0.0 &2.0$\pm$0.0 & 99.3$\pm$0.0 &99.5$\pm$0.0 &\underline{\bfseries 1} &65.3$\pm$0.0 &5.0$\pm$0.0 &99.1$\pm$0.0& 99.0$\pm$0.0 &\underline{\bfseries 1} \\
        {\bfseries ls1-GA} &65.0$\pm$1.9 &6432.1$\pm$6.4 & 99.1$\pm$0.2 &98.8$\pm$0.1 &\underline{\bfseries 1} &59.3$\pm$3.9 &6332.7$\pm$6.7& 99.0$\pm$0.0 &98.8$\pm$0.2 &\underline{\bfseries 1} \\
        {\bfseries ls1-EDA} &68.9$\pm$0.0 &6423.1$\pm$2.9 & 99.1$\pm$0.0 &99.4$\pm$0 &\underline{\bfseries 1} &61.35$\pm$0 &6323.2$\pm$3.2& 99.2$\pm$0.0 &99.1$\pm$0.0 &\underline{\bfseries 1} \\
        {\bfseries ls1-DDMLS(O)} &\underline{\bfseries61.5$\pm$0.0} &6417.8$\pm$16.3 & 99.2$\pm$0.0 &98.5$\pm$0 &\underline{\bfseries 1} &\underline{\bfseries53.9$\pm$2.3} &6318$\pm$65.9& 99.0$\pm$0.0 &98.6$\pm$0.1 &0.9 \\
        {\bfseries ls1-DDMLS(V1)} &63.4$\pm$1.4 &6417.8$\pm$16.3 & 99.6$\pm$0.0 &99.1$\pm$0.1 &\underline{\bfseries 1} &61.8$\pm$0.5 &6318$\pm$65.9& 99.7$\pm$0.3 &99.4$\pm$0.2 &\underline{\bfseries 1} \\
        {\bfseries ls1-DDMLS(V2)} &62.3$\pm$ 0.3 &6417.8$\pm$16.3 & 99.4$\pm$0.2 &98.9$\pm$0.3 &\underline{\bfseries 1} &54.8$\pm$2.8 &6318$\pm$65.9& 99.1$\pm$0.1 &98.8$\pm$0.1 &\underline{\bfseries 1} \\
        {\bfseries ls1-DDMLS(V3)} &62.9$\pm$0.8 &6403 & 99.2$\pm$0.1 &98.7$\pm$0.2 &\underline{\bfseries 1} &57.4$\pm$2.6 &6340$\pm$65.9& 99.1$\pm$0.0 &98.7$\pm$0.1 &\underline{\bfseries 1} \\
        
        \midrule
        {\bfseries ls2-Greedy} &85.7$\pm$0.0&3.0$\pm$0.0 & 99.9$\pm$0.0&99.9$\pm$0.0 &\underline{\bfseries 1} &85.7$\pm$0.0&3.0$\pm$0.0& 100.0$\pm$0.0&100.0$\pm$0.0  &\underline{\bfseries 1} \\
        {\bfseries ls2-GA} &76.9$\pm$2.7&22617.4$\pm$10.9 & 99.0$\pm$0.1&98.9$\pm$0.2 &\underline{\bfseries 1} &61.2$\pm$2.1&21327.1$\pm$20.3& 99.1$\pm$0.0&99.1$\pm$0.1  &\underline{\bfseries 1} \\
        {\bfseries ls2-EDA} &75.9$\pm$0.0&22604.3$\pm$3.3 & 99.0$\pm$0.0&98.9$\pm$0.0 &\underline{\bfseries 1} &59.2$\pm$0.3&21306.5$\pm$2.5& 99.2$\pm$0.0&99.2$\pm$0.1  &\underline{\bfseries 1} \\
        {\bfseries ls2-DDALS(O)} &\underline{\bfseries 71.9$\pm$0.3} &22591.3$\pm$63.6 & 99.0$\pm$0.0&99.0$\pm$0.0 &\underline{\bfseries 1} &\underline{\bfseries 57.2$\pm$0.0} &21233.5$\pm$51.9& 99.3$\pm$0.1&99.3$\pm$0.1  &\underline{\bfseries 1} \\
        {\bfseries ls2-DDALS(V1)} &76.6$\pm$1.1&22591.3$\pm$63.6 & 99.6$\pm$0.1&99.5$\pm$0.1 &\underline{\bfseries 1} &58.5$\pm$0.8&21233.5$\pm$51.9& 99.6$\pm$0.1&99.6$\pm$0.1  &\underline{\bfseries 1} \\
        {\bfseries ls2-DDALS(V2)} &72.9$\pm$0.6&22591.3$\pm$63.6 & 99.3$\pm$0.2&99.2$\pm$0.2 &\underline{\bfseries 1} &57.5$\pm$0.2&21233.5$\pm$51.9& 99.4$\pm$0.1&99.5$\pm$0.1  &\underline{\bfseries 1} \\
        {\bfseries ls2-DDALS(V3)} &73.1$\pm$0.7&22564 & 99.2$\pm$0.1&99.1$\pm$0.1 &\underline{\bfseries 1} &57.9$\pm$0.4&21245& 99.2$\pm$0.1&99.2$\pm$0.2  &\underline{\bfseries 1} \\
        
        \midrule
        {\bfseries ls3-Greedy} &158.3$\pm$0.0 &3.0$\pm$0.0 & 99.4$\pm$0.0 &99.3$\pm$0.0 &\underline{\bfseries 1} &158.3$\pm$0.0 &3.0$\pm$0.0& 99.6$\pm$0.0 &99.9$\pm$0.0 &\underline{\bfseries 1} \\
        {\bfseries ls3-GA} &119.4$\pm$3.9 &21008.5$\pm$7.5 & 99.1$\pm$0.1 &98.4$\pm$0.4 &0.7 &122.6$\pm$5.3 &21007.5$\pm$4.2 &99.0$\pm$0.1& 98.4$\pm$0.4   &0.5 \\
        {\bfseries ls3-EDA} &117.0$\pm$0.7 &21005.3$\pm$3.0 & 99.4$\pm$0.0 &98.9$\pm$0.0 &\underline{\bfseries 1} &115.0$\pm$2.0 &21004.8$\pm$2.6& 99.3$\pm$0.2 &99.0$\pm$0.1 &0.6 \\
        {\bfseries ls3-DDMLS(O)} &\underline{\bfseries112.0$\pm$1.2} &20942.2$\pm$211.2 & 99.1$\pm$0.0 &98.9$\pm$0.1 &0.9 &\underline{\bfseries111.4$\pm$ 1.5} &20956.2$\pm$193.7& 99.3$\pm$0.3 &98.8$\pm$0.4 &0.5 \\
        {\bfseries ls3-DDMLS(V1)} &119.0$\pm$2.7 &20942.2$\pm$211.2 & 99.5$\pm$0.1 &98.9$\pm$0.3 &0.9 &114.4$\pm$4.3 &20956.2$\pm$193.7& 99.5$\pm$0.2 &99.0$\pm$0.4 &0.7 \\
        {\bfseries ls3-DDMLS(V2)} &115.7$\pm$1.9 &20942.2$\pm$211.2 & 99.4$\pm$0.1 &98.8$\pm$0.1 &\underline{\bfseries 1} &111.8$\pm$ 1.7 &20956.2$\pm$193.7& 99.4$\pm$0.3 &98.9$\pm$0.4 &0.6 \\
        {\bfseries ls3-DDMLS(V3)} &115.6$\pm$1.5 &20810 & 99.2$\pm$0.1 &97.9$\pm$0.6 &0.3 &113.7$\pm$0.7 &20819& 99.0$\pm$0.0 &98.3$\pm$0.2 &0.1 \\
        
        \midrule
        {\bfseries ls4-Greedy} &207.3$\pm$0.0&13.0$\pm$0.0 & 99.0$\pm$0.0&98.9$\pm$0.0 &\underline{\bfseries 1} &195.3$\pm$0.0&17.0$\pm$0.0& 99.4$\pm$0.0&99.4$\pm$0.0  &\underline{\bfseries 1} \\
        {\bfseries ls4-GA} &208.5$\pm$8.3&48014.1$\pm$12.1 & 99.0$\pm$0.0&98.6$\pm$0.3&0.7&175.2$\pm$10.3&45014.6$\pm$10.2& 99.0$\pm$0.0&98.5$\pm$0.2  &0.6 \\
        {\bfseries ls4-EDA} &205.3$\pm$2.4&48004.5$\pm$3.3 & 99.0$\pm$0.0&99.0$\pm$0.0 &\underline{\bfseries 1} &163.8$\pm$1.3&45005.5$\pm$2.6& 99.1$\pm$0.0&98.8$\pm$0.0  &\underline{\bfseries 1} \\
        {\bfseries ls4-DDALS(O)} &\underline{\bfseries 194.9$\pm$0.8} &47962.8$\pm$67.7 & 99.0$\pm$0.0&98.9$\pm$0.0 &\underline{\bfseries 1} &\underline{\bfseries 163.3$\pm$3.6} &44822.1$\pm$583.5& 99.0$\pm$0.0&98.4$\pm$0.1  &0.3 \\
        {\bfseries ls4-DDALS(V1)} &203.2$\pm$2.0&47962.8$\pm$67.7 & 99.1$\pm$0.0&99.0$\pm$0.1 &\underline{\bfseries 1} &172.8$\pm$4.5&44822.1$\pm$583.5& 99.3$\pm$0.1&98.9$\pm$0.3  &0.8 \\
        {\bfseries ls4-DDALS(V2)} &196.5$\pm$1.3&47962.8$\pm$67.7 & 99.0$\pm$0.0&99.0$\pm$0.0 &\underline{\bfseries 1} &165.3$\pm$4.3&44822.1$\pm$583.5& 99.1$\pm$0.1&98.5$\pm$0.2  &0.3 \\
        {\bfseries ls4-DDALS(V3)} &196.9$\pm$2.1&48010 & 99.0$\pm$0.0&98.9$\pm$0.1 &\underline{\bfseries 1} &168.2$\pm$ 2.3&44174& 99.0$\pm$0.0& 98.9$\pm$ 0.1  &\underline{\bfseries 1} \\
        
        \midrule
        {\bfseries ls5-Greedy} &267.2$\pm$0.0 &11.0$\pm$0.0 & 99.1$\pm$0.0 &98.6$\pm$0.0 &\underline{\bfseries 1} &264.0$\pm$0.0 &12.0$\pm$0.0& 99.3$\pm$0.0 &98.9$\pm$0.0 &\underline{\bfseries 1} \\
        {\bfseries ls5-GA} &241.2$\pm$9.4 &81015.9$\pm$10.3 & 99.0$\pm$0.0 &98.6$\pm$0.1 &0.7 &219.9$\pm$10.5 &79012.6$\pm$7.5& 99.0$\pm$0.0 &98.5$\pm$0.2 &0.6 \\
        {\bfseries ls5-EDA} &220.3$\pm$0.3 &81002.4$\pm$2.4 & 99.1$\pm$0.0 &98.7$\pm$0.0 &\underline{\bfseries 1} &\underline{\bfseries198.7$\pm$1.8} &79005.3$\pm$3.1& 99.1$\pm$0.2 &98.6$\pm$0.2 &\underline{\bfseries 1}\\
        {\bfseries ls5-DDMLS(O)} &\underline{\bfseries218.2$\pm$1.0} &81144.5$\pm$283.4 & 99.0$\pm$0.0 &98.6$\pm$0.0 &\underline{\bfseries 1} &199.4$\pm$2.2 &78638.8$\pm$555.0& 99.0$\pm$0.1 &98.6$\pm$0.1 &0.9 \\
        {\bfseries ls5-DDMLS(V1)} &228.3$\pm$2.9 &81144.5$\pm$283.4 & 99.2$\pm$0.1 &98.9$\pm$0.2 &\underline{\bfseries 1} &209.3$\pm$5.8 &78638.8$\pm$555.0& 99.2$\pm$0.1 &98.9$\pm$0.2 &0.9 \\
        {\bfseries ls5-DDMLS(V2)} &220.3$\pm$3.0 &81144.5$\pm$283.4 & 99.1$\pm$0.1 &98.7$\pm$0.1 &\underline{\bfseries 1} &200.1$\pm$2.7 &78638.8$\pm$555.0& 99.1$\pm$0.1 &98.7$\pm$0.2 &\underline{\bfseries 1} \\
        {\bfseries ls5-DDMLS(V3)} &221.5$\pm$1.5 &81213 & 99.0$\pm$0.0 &98.7$\pm$0.0 &\underline{\bfseries 1} &202.2$\pm$3.1 &78784& 99.1$\pm$0.1 &98.9$\pm$0.2 &0.9 \\
        
        \midrule
        {\bfseries ls6-Greedy} &388.8$\pm$0.0 &4.0$\pm$0.0 & 99.1$\pm$0.0 &98.6$\pm$0.0 &\underline{\bfseries 1} &374.4$\pm$0.0 &12.0$\pm$0.0& 99.7$\pm$0.0 &99.6$\pm$0.0 &\underline{\bfseries 1} \\
        {\bfseries ls6-GA} &354.5$\pm$14.3 &125315.6$\pm$10.0 & 99.0$\pm$0.0 &98.2$\pm$0.1 &0 &325.3$\pm$10.5 &121508.9$\pm$4.8& 99.0$\pm$0.0 &98.2$\pm$0.2 &0.1 \\
        {\bfseries ls6-EDA} &337.8$\pm$0.9 &125305.4$\pm$2.4 & 99.0$\pm$0.0 &98.6$\pm$0.0 &\underline{\bfseries 1} &301.6$\pm$1.7 &121503.7$\pm$2.6& 99.1$\pm$0.1 &98.7$\pm$0.1 &\underline{\bfseries 1} \\
        {\bfseries ls6-DDMLS(O)} &\underline{\bfseries324.2$\pm$0.9} &125232.9$\pm$226.2 & 99.0$\pm$0.0 &98.1$\pm$0.0 &0 &\underline{\bfseries298.7$\pm$5.6} &121483.4$\pm$943.5& 99.0$\pm$0.0 &98.4$\pm$0.1 &0.3 \\
        {\bfseries ls6-DDMLS(V1)} &333.1$\pm$1.9 &125232.9$\pm$226.2 & 99.2$\pm$0.1 &98.6$\pm$0.2 &0.6 &311.5$\pm$7.0 &121483.4$\pm$943.5& 99.3$\pm$0.1 &98.6$\pm$0.3 &0.6 \\
        {\bfseries ls6-DDMLS(V2)} &325.5$\pm$2.7 &125232.9$\pm$226.2 & 99.0$\pm$0.1 &98.2$\pm$0.1 &0.1 &301.4$\pm$5.2 &121483.4$\pm$943.5& 99.1$\pm$0.2 &98.5$\pm$0.2 &0.5 \\
        {\bfseries ls6-DDMLS(V3)} &328.3$\pm$2.4 &124923 & 99.0$\pm$0.0 &98.3$\pm$0.1 &0 &304.7$\pm$3.2 &121648& 99.1$\pm$0.0 &98.4$\pm$0.1 &0.1 \\
        
        \bottomrule
        \end{tabular}
    }
    \label{fullresult:APP}
\end{table*}
% use section* for acknowledgment
% \section*{Acknowledgment}


% The authors would like to thank...


% Can use something like this to put references on a page
% by themselves when using endfloat and the captionsoff option.
\ifCLASSOPTIONcaptionsoff
  \newpage
\fi



% trigger a \newpage just before the given reference
% number - used to balance the columns on the last page
% adjust value as needed - may need to be readjusted if
% the document is modified later
%\IEEEtriggeratref{8}
% The "triggered" command can be changed if desired:
%\IEEEtriggercmd{\enlargethispage{-5in}}

% references section

\bibliographystyle{IEEEtran}
\bibliography{ref.bib}

% biography section
% 
% If you have an EPS/PDF photo (graphicx package needed) extra braces are
% needed around the contents of the optional argument to biography to prevent
% the LaTeX parser from getting confused when it sees the complicated
% \includegraphics command within an optional argument. (You could create
% your own custom macro containing the \includegraphics command to make things
% simpler here.)
%\begin{IEEEbiography}[{\includegraphics[width=1in,height=1.25in,clip,keepaspectratio]{mshell}}]{Michael Shell}
% or if you just want to reserve a space for a photo:

% \begin{IEEEbiography}{Xuanfeng Li}
% Biography text here.
% \end{IEEEbiography}

% % if you will not have a photo at all:
% \begin{IEEEbiographynophoto}{Shengcai Liu}
% Biography text here.
% \end{IEEEbiographynophoto}

% % insert where needed to balance the two columns on the last page with
% % biographies
% %\newpage

% \begin{IEEEbiographynophoto}{Ke Tang}
% Biography text here.
% \end{IEEEbiographynophoto}

% You can push biographies down or up by placing
% a \vfill before or after them. The appropriate
% use of \vfill depends on what kind of text is
% on the last page and whether or not the columns
% are being equalized.

%\vfill

% Can be used to pull up biographies so that the bottom of the last one
% is flush with the other column.
%\enlargethispage{-5in}



% that's all folks
\end{document}


